\section{孙吴：江面、继承与簿籍中的江东}
\label{sec:sun-wu-state-south}

\subsection{册封不能越过江面}

二二二年曹丕册封孙权为吴王，并要求质子，试图把长江下游政权纳入魏的
册命秩序。\eventanchor{TK-0222-WEI-INVESTS-SUN-QUAN}{魏册孙权}可见的不是
江东被洛阳接收，而是孙权借名义牵制蜀军、魏借名义等待江东让步的双向
试探。州郡吏员、水军与地方兵力并没有随一纸册命转移。孙权随后在二二九年
\eventanchor{TK-0229-SUN-QUAN-EMPEROR}{称帝}、\eventanchor{TK-0229-JIANYE-CAPITAL}{迁都建业}，
把下游交通、官僚与资源腹地置为
中枢；武昌仍保存面向上游的军事意义。两座城的并用，比“迁都完成”更能说明
江东必须在江面两端调度。

孙吴的可贵资源也是昂贵资源。长江让北军不能以陆路优势直接取胜，却要求
吴廷长期维持船只、港口、粮食、水手与沿岸城戍。二三〇年卫温、诸葛直
航海求夷洲、亶洲，许多军民病死；它显示远航的组织能力与海上动员的极限
同时存在。岛屿地望、航程和人员数不足以由后世地名或传记渲染裁定。

\subsection{继承将边防拖入朝廷}

孙登死后，孙和、孙霸各自形成宾客与官僚网络。二五〇年孙权废孙和、赐死
孙霸，改立孙亮。其结果可确认，宫廷言辞与各人罪责则多由后来的传记组织。
真正的制度难题在于：当成熟继承人和其支持者被清洗，江东必须把年幼皇帝的
军政连续性托付给辅政者，边防成败遂会立刻转换为朝廷人事的胜负。

二五二年诸葛恪在东兴抵御魏军，利用堤、防线和水陆地形守住濡须方向；
次年他围攻合肥新城，因疫病、补给与魏守备撤退。
\eventanchor{TK-0253-ZHUGE-KE-HEFEI}{合肥撤军}之后，诸葛恪被孙峻所杀。
这一串变化并非“江东善水战”或“权臣误国”可以解释：防守可依托既有节点，
远征则必须在季节、攻城器械、粮道和疾病中维持军队；失败还会改变谁有资格
控制禁军。二五八年孙綝废孙亮，\eventanchor{TK-0258-SUN-LIANG-DEPOSED}
{孙亮被废}，皇帝中心因而更加依赖不稳定的近侍与将领网络。

\subsection{江东并不只在建业}

走马楼吴简让人看到长沙临湘的户籍、田租和吏役；这些簿籍使江东国家能力
不再只由建业宫廷来代表。它们所能问的是谁被登记、何种米布钱和劳役被写入
官府程序，以及哪些家庭通过文书被看见；它们不能给出孙吴全境人口，也不能
代替山越、流民、沿海人群的声音。地方行政正是在这种不均匀的可见性中运行。

二六四年孙皓即位，西陵、荆州和武昌的防务仍能使一名边将、一处峡口左右
战略。二七二年陆抗收复西陵，表明吴的水陆网络未立即崩散；但继承清洗、
上游防务和中央调度长期互相消耗。它为西晋后来分路沿江推进留下了条件，
却不能倒过来把孙吴末年每一项地方制度写成完全失效。

\begin{sourcenote}
\sourcebook{三国志·吴书}提供孙吴称帝、继承与军事行动的主干，却由以魏为
框架的编纂保存。走马楼吴简能校正“只有宫廷”的视野，但其地域和机构很窄；
城址与航路考古能提示空间条件，不能补成一张精确疆界或户口表。孙权称帝的
证据范围见本卷保留的 `TK-0229-SUN-QUAN-EMPEROR` 证据护照。
\end{sourcenote}

% DEPTH-EXPANSION-2026
\section{孙吴：江面把政权连成网络，也把它拉开：把事件簇读成可持续的约束}

这一节不是把账本改写成一排胜负。它把已登记的事件放回长江上下游、淮南接口、荆州峡口和长沙文书所见的地方，追问资源如何被调用、选择如何受限、制度怎样在成功与撤退之间留下收益和亏损。每一个节点仍以账本的日期与材料为界；扩展的目的，是让读者在一处命令、一段行程或一份文书中同时看见空间、行动者与后来必须承接的压力。


\subsection{把曹丕册封孙权为吴王，试图以册命、质子放回资源、道路与制度}
\label{reading:tk-0222-wei-invests-sun-quan}

黄初三年（222年），账本所记的\eventref{TK-0222-WEI-INVESTS-SUN-QUAN}{曹丕册封孙权为吴王，试图以册命、质子要求和荆州牧名义约束……}，首先不是抽象地图上移动的一枚棋子。读者应先把目光放在建业、武昌、濡须与长江港口：在这里，一道命令若要抵达，就必须穿过长江上下游、淮南接口、荆州峡口和长沙文书所见的地方。曹丕册封孙权为吴王，试图以册命、质子要求和荆州牧名义约束长江下游政权。这一结果给出可核验的时间骨架；它同时提出更难的问题：谁能调动物资与人手，谁只能等待别人的决定，谁又在道路、水面或官署的缝隙里重新计算自己的位置。

从结构上说，舟师、港口、沿江城戍、州郡簿籍和能够随水路移动的粮食并不是储存在国库里便可自动使用的力量。它们需要在不同地点被登记、运出、交接、守护和补充；其中任何一环延误，远方的军队或使节就会面对与中枢不同的现实。魏在刘备东征时希望利用孙刘冲突，以名义封授换取孙权承认北方皇帝和交出人质。这段前因不应被理解为单一人物预先设计的剧本，而是一组已经存在的条件：空间给出距离，制度分配发令权，地方社会承担把命令做成行动的成本。

行动者面对的选择因而很具体。对于这一政权或集团而言，关键不是笼统地“进”或“退”，而是把下游中枢、上游防线与继承安排维持在同一条可运转的江面上。若冒险推进，可能赢得一处节点、一个名义或一次短期主动；若保留力量，又可能把接口让给竞争者，或让既有同盟与部属怀疑中枢的能力。水军维持、岸防、疾病、远征补给和宫廷清洗互相放大不是事后的修辞，它在当时已构成选择的价格，只是材料很少留下所有人的意见与犹疑。

在这一限制下，曹丕册封孙权为吴王，试图以册命、质子要求和荆州牧名义约束长江下游政权。把事实写到这里仍须克制：正史可以确认的通常是任命、出军、围守、退还、死亡、册封或政权转换等骨架；营中每一段传话、每一份账目和每一处战阵，往往没有同样坚实的记录。真正有解释力的，并非补足一个好看的传奇，而是观察这一骨架怎样被长江上下游、淮南接口、荆州峡口和长沙文书所见的地方的摩擦推向已经发生的结果。

直接效果首先落在可见的节点上。孙权暂借魏的封号牵制蜀汉而未交出统治权；魏吴间的宗藩语言很快让位于边境战争。一个城、一个渡口、一个官署、一名辅政者或一条通道的状态改变，会立刻重排相邻行动者的风险：前线可能得到喘息，近畿可能获得声望，地方官却要重新面对征发、守备或接收的命令。事件的“结束”因此只是中央记事的一条句号；在运输、任官和地方关系中，它刚刚开始产生后果。

中期影响并不等于结果的线性放大。成功会使中枢相信某条路线、某位将领或某种制度安排值得继续投资，也可能把资源更加集中到已经拥挤的方向；失利则可能迫使它缩短战线、改写同盟或重排辅政权力。就此而言，禁军、辅政、太子网络、州郡与水陆指挥的连接既是此次行动的条件，也是它留下的制度遗产。国家能力并不是“有”或“无”，而是在每一次调用之后增添某种可重复的程序，同时制造另一处更难处理的依赖。

把镜头拉近，船只、粮米、劳役和地方登记使江面上的战略最终落到家庭与县廷。史书多从朝廷、将帅和使节写起，难以连续记录沿线家庭如何应付粮食、役使、迁徙和不确定的安全；这并不意味着他们只是背景。恰恰是这些无法被完整统计的劳动，使一条道路、一支舟师或一座城戍能够在若干季节内维持。读者不宜凭缺失的户口表虚构精确的数字，却应把这种材料缺口本身看作国家抵达社会并不均匀的证据。

这一事件也改变了与他者相处的方式。把下游中枢、上游防线与继承安排维持在同一条可运转的江面上在边境上经常意味着对方也必须改算自己的粮道、使节、驻军与继承风险。即使双方暂时不战，册命、招纳、投降或撤军也会传递有关实力与脆弱性的信息。外交在这里不是战争之外的礼节，而是以较低成本测试关系的一种行动；战争则是在谈判、承诺与防守网被证明不足之后，迫使对手承担更高成本的手段。

制度所得必须和制度损失一起衡量。此次行动若使命令链更清楚、边防节点更稳固、任官更能配合资源，它就扩大了中枢的可用能力；但若它迫使权力过度依赖少数近侍、边将、转输官或地方合作者，也会让下一次继承、叛乱或远征更难收束。水军维持、岸防、疾病、远征补给和宫廷清洗互相放大由此不只是军费或伤亡的代名词，更是信息集中、地方弹性与政治信任被逐步消耗的过程。

材料在这里决定叙述能走多远。本段的基本年序和行动依据是《三国志·吴主传》；《三国志·文帝纪》；《资治通鉴》魏纪。它们能够相互校验一些结果，却不是互不相干的现场档案：纪传有编纂目的，后出的通鉴重组了时间，地方文献与考古又只照亮有限地点。因此，《三国志》吴书保存宫廷主线，走马楼吴简和区域考古仅能从局部校正国家叙事。引用多部书不自动制造全知视角；不同材料的沉默、重叠与矛盾，正是解释的一部分。

争议尤其集中在动机、数字与细节。地望、兵额、海上航程、宫廷对话和江东全境的户口都不能由一项材料裁定。较稳妥的写法是把可确认的结果与解释性的推断分开：可以说行动者在既有资源和压力下作出某种选择，也必须承认我们通常无法知道每个人在当时究竟相信什么、害怕什么。这样的保留并不会削弱历史，反而使读者看见选择为何有限、失败为何未必等于无能、胜利为何不必然证明先见。

如果把这一节点接回全章，它所显示的是一条反馈回路：舟师、港口、沿江城戍、州郡簿籍和能够随水路移动的粮食使政权能够作出行动；行动又把压力送回禁军、辅政、太子网络、州郡与水陆指挥的连接和船只、粮米、劳役和地方登记使江面上的战略最终落到家庭与县廷；新形成的权力关系反过来决定下一次资源由谁调用、沿哪条路线流动。三国史最值得阅读之处，正在于这种循环不断把宏大的名号压回具体的河流、山口、港口、文书与人事之中。曹丕册封孙权为吴王，试图以册命、质子要求和荆州牧名义约束长江下游政权不是孤立条目，而是这条回路在黄初三年（222年）留下的一处可见折痕。


\subsection{从“陆逊在石亭击败曹休所率魏军，魏的江淮攻势受挫”看见行动的边界}
\label{reading:tk-0228-shiting}

太和二年（228年），账本所记的\eventanchor{TK-0228-SHITING}{陆逊在石亭击败曹休所率魏军，魏的江淮攻势受挫}，首先不是抽象地图上移动的一枚棋子。读者应先把目光放在建业、武昌、濡须与长江港口：在这里，一道命令若要抵达，就必须穿过长江上下游、淮南接口、荆州峡口和长沙文书所见的地方。陆逊在石亭击败曹休所率魏军，魏的江淮攻势受挫。这一结果给出可核验的时间骨架；它同时提出更难的问题：谁能调动物资与人手，谁只能等待别人的决定，谁又在道路、水面或官署的缝隙里重新计算自己的位置。

从结构上说，舟师、港口、沿江城戍、州郡簿籍和能够随水路移动的粮食并不是储存在国库里便可自动使用的力量。它们需要在不同地点被登记、运出、交接、守护和补充；其中任何一环延误，远方的军队或使节就会面对与中枢不同的现实。魏试图沿江淮推进并利用吴方内部关系，孙吴以陆逊、朱桓等将领掌握水陆节点和情报优势。这段前因不应被理解为单一人物预先设计的剧本，而是一组已经存在的条件：空间给出距离，制度分配发令权，地方社会承担把命令做成行动的成本。

行动者面对的选择因而很具体。对于这一政权或集团而言，关键不是笼统地“进”或“退”，而是把下游中枢、上游防线与继承安排维持在同一条可运转的江面上。若冒险推进，可能赢得一处节点、一个名义或一次短期主动；若保留力量，又可能把接口让给竞争者，或让既有同盟与部属怀疑中枢的能力。水军维持、岸防、疾病、远征补给和宫廷清洗互相放大不是事后的修辞，它在当时已构成选择的价格，只是材料很少留下所有人的意见与犹疑。

在这一限制下，陆逊在石亭击败曹休所率魏军，魏的江淮攻势受挫。把事实写到这里仍须克制：正史可以确认的通常是任命、出军、围守、退还、死亡、册封或政权转换等骨架；营中每一段传话、每一份账目和每一处战阵，往往没有同样坚实的记录。真正有解释力的，并非补足一个好看的传奇，而是观察这一骨架怎样被长江上下游、淮南接口、荆州峡口和长沙文书所见的地方的摩擦推向已经发生的结果。

直接效果首先落在可见的节点上。魏军后撤并调整淮南防御，孙吴证明其可在江北作战；诈降情节主要见于传记叙事，须区分战果与戏剧化细节。一个城、一个渡口、一个官署、一名辅政者或一条通道的状态改变，会立刻重排相邻行动者的风险：前线可能得到喘息，近畿可能获得声望，地方官却要重新面对征发、守备或接收的命令。事件的“结束”因此只是中央记事的一条句号；在运输、任官和地方关系中，它刚刚开始产生后果。

中期影响并不等于结果的线性放大。成功会使中枢相信某条路线、某位将领或某种制度安排值得继续投资，也可能把资源更加集中到已经拥挤的方向；失利则可能迫使它缩短战线、改写同盟或重排辅政权力。就此而言，禁军、辅政、太子网络、州郡与水陆指挥的连接既是此次行动的条件，也是它留下的制度遗产。国家能力并不是“有”或“无”，而是在每一次调用之后增添某种可重复的程序，同时制造另一处更难处理的依赖。

把镜头拉近，船只、粮米、劳役和地方登记使江面上的战略最终落到家庭与县廷。史书多从朝廷、将帅和使节写起，难以连续记录沿线家庭如何应付粮食、役使、迁徙和不确定的安全；这并不意味着他们只是背景。恰恰是这些无法被完整统计的劳动，使一条道路、一支舟师或一座城戍能够在若干季节内维持。读者不宜凭缺失的户口表虚构精确的数字，却应把这种材料缺口本身看作国家抵达社会并不均匀的证据。

这一事件也改变了与他者相处的方式。把下游中枢、上游防线与继承安排维持在同一条可运转的江面上在边境上经常意味着对方也必须改算自己的粮道、使节、驻军与继承风险。即使双方暂时不战，册命、招纳、投降或撤军也会传递有关实力与脆弱性的信息。外交在这里不是战争之外的礼节，而是以较低成本测试关系的一种行动；战争则是在谈判、承诺与防守网被证明不足之后，迫使对手承担更高成本的手段。

制度所得必须和制度损失一起衡量。此次行动若使命令链更清楚、边防节点更稳固、任官更能配合资源，它就扩大了中枢的可用能力；但若它迫使权力过度依赖少数近侍、边将、转输官或地方合作者，也会让下一次继承、叛乱或远征更难收束。水军维持、岸防、疾病、远征补给和宫廷清洗互相放大由此不只是军费或伤亡的代名词，更是信息集中、地方弹性与政治信任被逐步消耗的过程。

材料在这里决定叙述能走多远。本段的基本年序和行动依据是《三国志·吴主传》；《三国志·陆逊传》；《三国志·曹休传》；《资治通鉴》魏纪。它们能够相互校验一些结果，却不是互不相干的现场档案：纪传有编纂目的，后出的通鉴重组了时间，地方文献与考古又只照亮有限地点。因此，《三国志》吴书保存宫廷主线，走马楼吴简和区域考古仅能从局部校正国家叙事。引用多部书不自动制造全知视角；不同材料的沉默、重叠与矛盾，正是解释的一部分。

争议尤其集中在动机、数字与细节。地望、兵额、海上航程、宫廷对话和江东全境的户口都不能由一项材料裁定。较稳妥的写法是把可确认的结果与解释性的推断分开：可以说行动者在既有资源和压力下作出某种选择，也必须承认我们通常无法知道每个人在当时究竟相信什么、害怕什么。这样的保留并不会削弱历史，反而使读者看见选择为何有限、失败为何未必等于无能、胜利为何不必然证明先见。

如果把这一节点接回全章，它所显示的是一条反馈回路：舟师、港口、沿江城戍、州郡簿籍和能够随水路移动的粮食使政权能够作出行动；行动又把压力送回禁军、辅政、太子网络、州郡与水陆指挥的连接和船只、粮米、劳役和地方登记使江面上的战略最终落到家庭与县廷；新形成的权力关系反过来决定下一次资源由谁调用、沿哪条路线流动。三国史最值得阅读之处，正在于这种循环不断把宏大的名号压回具体的河流、山口、港口、文书与人事之中。陆逊在石亭击败曹休所率魏军，魏的江淮攻势受挫不是孤立条目，而是这条回路在太和二年（228年）留下的一处可见折痕。


\subsection{建业、武昌、濡须与长江港口：一次转折的可达与不可达}
\label{reading:tk-0229-sun-quan-emperor}

黄龙元年四月（229年），账本所记的\eventref{TK-0229-SUN-QUAN-EMPEROR}{孙权在武昌即皇帝位，建元黄龙，孙吴正式提出独立帝国名号}，首先不是抽象地图上移动的一枚棋子。读者应先把目光放在建业、武昌、濡须与长江港口：在这里，一道命令若要抵达，就必须穿过长江上下游、淮南接口、荆州峡口和长沙文书所见的地方。孙权在武昌即皇帝位，建元黄龙，孙吴正式提出独立帝国名号。这一结果给出可核验的时间骨架；它同时提出更难的问题：谁能调动物资与人手，谁只能等待别人的决定，谁又在道路、水面或官署的缝隙里重新计算自己的位置。

从结构上说，舟师、港口、沿江城戍、州郡簿籍和能够随水路移动的粮食并不是储存在国库里便可自动使用的力量。它们需要在不同地点被登记、运出、交接、守护和补充；其中任何一环延误，远方的军队或使节就会面对与中枢不同的现实。魏的册封与质子要求不能满足孙权的自主统治，蜀汉亦已称帝，江东政权需要以独立名号配置官爵和外交。这段前因不应被理解为单一人物预先设计的剧本，而是一组已经存在的条件：空间给出距离，制度分配发令权，地方社会承担把命令做成行动的成本。

行动者面对的选择因而很具体。对于这一政权或集团而言，关键不是笼统地“进”或“退”，而是把下游中枢、上游防线与继承安排维持在同一条可运转的江面上。若冒险推进，可能赢得一处节点、一个名义或一次短期主动；若保留力量，又可能把接口让给竞争者，或让既有同盟与部属怀疑中枢的能力。水军维持、岸防、疾病、远征补给和宫廷清洗互相放大不是事后的修辞，它在当时已构成选择的价格，只是材料很少留下所有人的意见与犹疑。

在这一限制下，孙权在武昌即皇帝位，建元黄龙，孙吴正式提出独立帝国名号。把事实写到这里仍须克制：正史可以确认的通常是任命、出军、围守、退还、死亡、册封或政权转换等骨架；营中每一段传话、每一份账目和每一处战阵，往往没有同样坚实的记录。真正有解释力的，并非补足一个好看的传奇，而是观察这一骨架怎样被长江上下游、淮南接口、荆州峡口和长沙文书所见的地方的摩擦推向已经发生的结果。

直接效果首先落在可见的节点上。孙吴与魏、蜀汉形成三个皇帝中心；称帝强化建国仪式，却未消除长江中上游与淮南的军事压力。一个城、一个渡口、一个官署、一名辅政者或一条通道的状态改变，会立刻重排相邻行动者的风险：前线可能得到喘息，近畿可能获得声望，地方官却要重新面对征发、守备或接收的命令。事件的“结束”因此只是中央记事的一条句号；在运输、任官和地方关系中，它刚刚开始产生后果。

中期影响并不等于结果的线性放大。成功会使中枢相信某条路线、某位将领或某种制度安排值得继续投资，也可能把资源更加集中到已经拥挤的方向；失利则可能迫使它缩短战线、改写同盟或重排辅政权力。就此而言，禁军、辅政、太子网络、州郡与水陆指挥的连接既是此次行动的条件，也是它留下的制度遗产。国家能力并不是“有”或“无”，而是在每一次调用之后增添某种可重复的程序，同时制造另一处更难处理的依赖。

把镜头拉近，船只、粮米、劳役和地方登记使江面上的战略最终落到家庭与县廷。史书多从朝廷、将帅和使节写起，难以连续记录沿线家庭如何应付粮食、役使、迁徙和不确定的安全；这并不意味着他们只是背景。恰恰是这些无法被完整统计的劳动，使一条道路、一支舟师或一座城戍能够在若干季节内维持。读者不宜凭缺失的户口表虚构精确的数字，却应把这种材料缺口本身看作国家抵达社会并不均匀的证据。

这一事件也改变了与他者相处的方式。把下游中枢、上游防线与继承安排维持在同一条可运转的江面上在边境上经常意味着对方也必须改算自己的粮道、使节、驻军与继承风险。即使双方暂时不战，册命、招纳、投降或撤军也会传递有关实力与脆弱性的信息。外交在这里不是战争之外的礼节，而是以较低成本测试关系的一种行动；战争则是在谈判、承诺与防守网被证明不足之后，迫使对手承担更高成本的手段。

制度所得必须和制度损失一起衡量。此次行动若使命令链更清楚、边防节点更稳固、任官更能配合资源，它就扩大了中枢的可用能力；但若它迫使权力过度依赖少数近侍、边将、转输官或地方合作者，也会让下一次继承、叛乱或远征更难收束。水军维持、岸防、疾病、远征补给和宫廷清洗互相放大由此不只是军费或伤亡的代名词，更是信息集中、地方弹性与政治信任被逐步消耗的过程。

材料在这里决定叙述能走多远。本段的基本年序和行动依据是《三国志·吴主传》；《三国志·孙登传》；《资治通鉴》魏纪。它们能够相互校验一些结果，却不是互不相干的现场档案：纪传有编纂目的，后出的通鉴重组了时间，地方文献与考古又只照亮有限地点。因此，《三国志》吴书保存宫廷主线，走马楼吴简和区域考古仅能从局部校正国家叙事。引用多部书不自动制造全知视角；不同材料的沉默、重叠与矛盾，正是解释的一部分。

争议尤其集中在动机、数字与细节。地望、兵额、海上航程、宫廷对话和江东全境的户口都不能由一项材料裁定。较稳妥的写法是把可确认的结果与解释性的推断分开：可以说行动者在既有资源和压力下作出某种选择，也必须承认我们通常无法知道每个人在当时究竟相信什么、害怕什么。这样的保留并不会削弱历史，反而使读者看见选择为何有限、失败为何未必等于无能、胜利为何不必然证明先见。

如果把这一节点接回全章，它所显示的是一条反馈回路：舟师、港口、沿江城戍、州郡簿籍和能够随水路移动的粮食使政权能够作出行动；行动又把压力送回禁军、辅政、太子网络、州郡与水陆指挥的连接和船只、粮米、劳役和地方登记使江面上的战略最终落到家庭与县廷；新形成的权力关系反过来决定下一次资源由谁调用、沿哪条路线流动。三国史最值得阅读之处，正在于这种循环不断把宏大的名号压回具体的河流、山口、港口、文书与人事之中。孙权在武昌即皇帝位，建元黄龙，孙吴正式提出独立帝国名号不是孤立条目，而是这条回路在黄龙元年四月（229年）留下的一处可见折痕。


\subsection{胜负之外：禁军、辅政、太子网络、州郡与水陆指挥的连接的压力}
\label{reading:tk-0229-jianye-capital}

黄龙元年（229年），账本所记的\eventref{TK-0229-JIANYE-CAPITAL}{孙权将都城由武昌迁回建业，长江下游成为孙吴中枢}，首先不是抽象地图上移动的一枚棋子。读者应先把目光放在建业、武昌、濡须与长江港口：在这里，一道命令若要抵达，就必须穿过长江上下游、淮南接口、荆州峡口和长沙文书所见的地方。孙权将都城由武昌迁回建业，长江下游成为孙吴中枢。这一结果给出可核验的时间骨架；它同时提出更难的问题：谁能调动物资与人手，谁只能等待别人的决定，谁又在道路、水面或官署的缝隙里重新计算自己的位置。

从结构上说，舟师、港口、沿江城戍、州郡簿籍和能够随水路移动的粮食并不是储存在国库里便可自动使用的力量。它们需要在不同地点被登记、运出、交接、守护和补充；其中任何一环延误，远方的军队或使节就会面对与中枢不同的现实。孙吴称帝后须在江东腹地建立较稳定的官僚与交通中心，同时保留武昌对上游的军事作用。这段前因不应被理解为单一人物预先设计的剧本，而是一组已经存在的条件：空间给出距离，制度分配发令权，地方社会承担把命令做成行动的成本。

行动者面对的选择因而很具体。对于这一政权或集团而言，关键不是笼统地“进”或“退”，而是把下游中枢、上游防线与继承安排维持在同一条可运转的江面上。若冒险推进，可能赢得一处节点、一个名义或一次短期主动；若保留力量，又可能把接口让给竞争者，或让既有同盟与部属怀疑中枢的能力。水军维持、岸防、疾病、远征补给和宫廷清洗互相放大不是事后的修辞，它在当时已构成选择的价格，只是材料很少留下所有人的意见与犹疑。

在这一限制下，孙权将都城由武昌迁回建业，长江下游成为孙吴中枢。把事实写到这里仍须克制：正史可以确认的通常是任命、出军、围守、退还、死亡、册封或政权转换等骨架；营中每一段传话、每一份账目和每一处战阵，往往没有同样坚实的记录。真正有解释力的，并非补足一个好看的传奇，而是观察这一骨架怎样被长江上下游、淮南接口、荆州峡口和长沙文书所见的地方的摩擦推向已经发生的结果。

直接效果首先落在可见的节点上。建业成为孙吴及后世建康政治传统的起点；城址材料可校验城市演变，不能直接复原全部三国宫城布局。一个城、一个渡口、一个官署、一名辅政者或一条通道的状态改变，会立刻重排相邻行动者的风险：前线可能得到喘息，近畿可能获得声望，地方官却要重新面对征发、守备或接收的命令。事件的“结束”因此只是中央记事的一条句号；在运输、任官和地方关系中，它刚刚开始产生后果。

中期影响并不等于结果的线性放大。成功会使中枢相信某条路线、某位将领或某种制度安排值得继续投资，也可能把资源更加集中到已经拥挤的方向；失利则可能迫使它缩短战线、改写同盟或重排辅政权力。就此而言，禁军、辅政、太子网络、州郡与水陆指挥的连接既是此次行动的条件，也是它留下的制度遗产。国家能力并不是“有”或“无”，而是在每一次调用之后增添某种可重复的程序，同时制造另一处更难处理的依赖。

把镜头拉近，船只、粮米、劳役和地方登记使江面上的战略最终落到家庭与县廷。史书多从朝廷、将帅和使节写起，难以连续记录沿线家庭如何应付粮食、役使、迁徙和不确定的安全；这并不意味着他们只是背景。恰恰是这些无法被完整统计的劳动，使一条道路、一支舟师或一座城戍能够在若干季节内维持。读者不宜凭缺失的户口表虚构精确的数字，却应把这种材料缺口本身看作国家抵达社会并不均匀的证据。

这一事件也改变了与他者相处的方式。把下游中枢、上游防线与继承安排维持在同一条可运转的江面上在边境上经常意味着对方也必须改算自己的粮道、使节、驻军与继承风险。即使双方暂时不战，册命、招纳、投降或撤军也会传递有关实力与脆弱性的信息。外交在这里不是战争之外的礼节，而是以较低成本测试关系的一种行动；战争则是在谈判、承诺与防守网被证明不足之后，迫使对手承担更高成本的手段。

制度所得必须和制度损失一起衡量。此次行动若使命令链更清楚、边防节点更稳固、任官更能配合资源，它就扩大了中枢的可用能力；但若它迫使权力过度依赖少数近侍、边将、转输官或地方合作者，也会让下一次继承、叛乱或远征更难收束。水军维持、岸防、疾病、远征补给和宫廷清洗互相放大由此不只是军费或伤亡的代名词，更是信息集中、地方弹性与政治信任被逐步消耗的过程。

材料在这里决定叙述能走多远。本段的基本年序和行动依据是《三国志·吴主传》；《建康实录》；南京地区六朝城址考古材料。它们能够相互校验一些结果，却不是互不相干的现场档案：纪传有编纂目的，后出的通鉴重组了时间，地方文献与考古又只照亮有限地点。因此，《三国志》吴书保存宫廷主线，走马楼吴简和区域考古仅能从局部校正国家叙事。引用多部书不自动制造全知视角；不同材料的沉默、重叠与矛盾，正是解释的一部分。

争议尤其集中在动机、数字与细节。地望、兵额、海上航程、宫廷对话和江东全境的户口都不能由一项材料裁定。较稳妥的写法是把可确认的结果与解释性的推断分开：可以说行动者在既有资源和压力下作出某种选择，也必须承认我们通常无法知道每个人在当时究竟相信什么、害怕什么。这样的保留并不会削弱历史，反而使读者看见选择为何有限、失败为何未必等于无能、胜利为何不必然证明先见。

如果把这一节点接回全章，它所显示的是一条反馈回路：舟师、港口、沿江城戍、州郡簿籍和能够随水路移动的粮食使政权能够作出行动；行动又把压力送回禁军、辅政、太子网络、州郡与水陆指挥的连接和船只、粮米、劳役和地方登记使江面上的战略最终落到家庭与县廷；新形成的权力关系反过来决定下一次资源由谁调用、沿哪条路线流动。三国史最值得阅读之处，正在于这种循环不断把宏大的名号压回具体的河流、山口、港口、文书与人事之中。孙权将都城由武昌迁回建业，长江下游成为孙吴中枢不是孤立条目，而是这条回路在黄龙元年（229年）留下的一处可见折痕。

\begin{sourcenote}
这一组事件的来源路径并不相同。《三国志》吴书保存宫廷主线，走马楼吴简和区域考古仅能从局部校正国家叙事。阅读时应先区分能确认的年序和结果，与只能作为解释线索的动机、数量及局部过程；这比用一个完整无缝的故事覆盖材料间隙更能说明国家能力如何实际运作。
\end{sourcenote}


\subsection{从“孙权遣卫温、诸葛直率舟师求夷洲、亶洲，航海军”看见行动的边界}
\label{reading:tk-0230-wei-wen-maritime}

黄龙二年（230年），账本所记的\eventanchor{TK-0230-WEI-WEN-MARITIME}{孙权遣卫温、诸葛直率舟师求夷洲、亶洲，航海军民多病死而未……}，首先不是抽象地图上移动的一枚棋子。读者应先把目光放在建业、武昌、濡须与长江港口：在这里，一道命令若要抵达，就必须穿过长江上下游、淮南接口、荆州峡口和长沙文书所见的地方。孙权遣卫温、诸葛直率舟师求夷洲、亶洲，航海军民多病死而未获预期成果。这一结果给出可核验的时间骨架；它同时提出更难的问题：谁能调动物资与人手，谁只能等待别人的决定，谁又在道路、水面或官署的缝隙里重新计算自己的位置。

从结构上说，舟师、港口、沿江城戍、州郡簿籍和能够随水路移动的粮食并不是储存在国库里便可自动使用的力量。它们需要在不同地点被登记、运出、交接、守护和补充；其中任何一环延误，远方的军队或使节就会面对与中枢不同的现实。孙吴拥有沿海舟师并寻求扩展人力、资源和外部信息，但对外海环境与岛屿社会缺乏稳定补给条件。这段前因不应被理解为单一人物预先设计的剧本，而是一组已经存在的条件：空间给出距离，制度分配发令权，地方社会承担把命令做成行动的成本。

行动者面对的选择因而很具体。对于这一政权或集团而言，关键不是笼统地“进”或“退”，而是把下游中枢、上游防线与继承安排维持在同一条可运转的江面上。若冒险推进，可能赢得一处节点、一个名义或一次短期主动；若保留力量，又可能把接口让给竞争者，或让既有同盟与部属怀疑中枢的能力。水军维持、岸防、疾病、远征补给和宫廷清洗互相放大不是事后的修辞，它在当时已构成选择的价格，只是材料很少留下所有人的意见与犹疑。

在这一限制下，孙权遣卫温、诸葛直率舟师求夷洲、亶洲，航海军民多病死而未获预期成果。把事实写到这里仍须克制：正史可以确认的通常是任命、出军、围守、退还、死亡、册封或政权转换等骨架；营中每一段传话、每一份账目和每一处战阵，往往没有同样坚实的记录。真正有解释力的，并非补足一个好看的传奇，而是观察这一骨架怎样被长江上下游、淮南接口、荆州峡口和长沙文书所见的地方的摩擦推向已经发生的结果。

直接效果首先落在可见的节点上。远航显示孙吴政权的海上动员能力及其限度；夷洲是否对应今台湾等问题必须保留文本地理的不确定性。一个城、一个渡口、一个官署、一名辅政者或一条通道的状态改变，会立刻重排相邻行动者的风险：前线可能得到喘息，近畿可能获得声望，地方官却要重新面对征发、守备或接收的命令。事件的“结束”因此只是中央记事的一条句号；在运输、任官和地方关系中，它刚刚开始产生后果。

中期影响并不等于结果的线性放大。成功会使中枢相信某条路线、某位将领或某种制度安排值得继续投资，也可能把资源更加集中到已经拥挤的方向；失利则可能迫使它缩短战线、改写同盟或重排辅政权力。就此而言，禁军、辅政、太子网络、州郡与水陆指挥的连接既是此次行动的条件，也是它留下的制度遗产。国家能力并不是“有”或“无”，而是在每一次调用之后增添某种可重复的程序，同时制造另一处更难处理的依赖。

把镜头拉近，船只、粮米、劳役和地方登记使江面上的战略最终落到家庭与县廷。史书多从朝廷、将帅和使节写起，难以连续记录沿线家庭如何应付粮食、役使、迁徙和不确定的安全；这并不意味着他们只是背景。恰恰是这些无法被完整统计的劳动，使一条道路、一支舟师或一座城戍能够在若干季节内维持。读者不宜凭缺失的户口表虚构精确的数字，却应把这种材料缺口本身看作国家抵达社会并不均匀的证据。

这一事件也改变了与他者相处的方式。把下游中枢、上游防线与继承安排维持在同一条可运转的江面上在边境上经常意味着对方也必须改算自己的粮道、使节、驻军与继承风险。即使双方暂时不战，册命、招纳、投降或撤军也会传递有关实力与脆弱性的信息。外交在这里不是战争之外的礼节，而是以较低成本测试关系的一种行动；战争则是在谈判、承诺与防守网被证明不足之后，迫使对手承担更高成本的手段。

制度所得必须和制度损失一起衡量。此次行动若使命令链更清楚、边防节点更稳固、任官更能配合资源，它就扩大了中枢的可用能力；但若它迫使权力过度依赖少数近侍、边将、转输官或地方合作者，也会让下一次继承、叛乱或远征更难收束。水军维持、岸防、疾病、远征补给和宫廷清洗互相放大由此不只是军费或伤亡的代名词，更是信息集中、地方弹性与政治信任被逐步消耗的过程。

材料在这里决定叙述能走多远。本段的基本年序和行动依据是《三国志·吴主传》；《三国志·吴主传》裴松之注引《临海水土志》；《资治通鉴》魏纪。它们能够相互校验一些结果，却不是互不相干的现场档案：纪传有编纂目的，后出的通鉴重组了时间，地方文献与考古又只照亮有限地点。因此，《三国志》吴书保存宫廷主线，走马楼吴简和区域考古仅能从局部校正国家叙事。引用多部书不自动制造全知视角；不同材料的沉默、重叠与矛盾，正是解释的一部分。

争议尤其集中在动机、数字与细节。地望、兵额、海上航程、宫廷对话和江东全境的户口都不能由一项材料裁定。较稳妥的写法是把可确认的结果与解释性的推断分开：可以说行动者在既有资源和压力下作出某种选择，也必须承认我们通常无法知道每个人在当时究竟相信什么、害怕什么。这样的保留并不会削弱历史，反而使读者看见选择为何有限、失败为何未必等于无能、胜利为何不必然证明先见。

如果把这一节点接回全章，它所显示的是一条反馈回路：舟师、港口、沿江城戍、州郡簿籍和能够随水路移动的粮食使政权能够作出行动；行动又把压力送回禁军、辅政、太子网络、州郡与水陆指挥的连接和船只、粮米、劳役和地方登记使江面上的战略最终落到家庭与县廷；新形成的权力关系反过来决定下一次资源由谁调用、沿哪条路线流动。三国史最值得阅读之处，正在于这种循环不断把宏大的名号压回具体的河流、山口、港口、文书与人事之中。孙权遣卫温、诸葛直率舟师求夷洲、亶洲，航海军民多病死而未获预期成果不是孤立条目，而是这条回路在黄龙二年（230年）留下的一处可见折痕。


\subsection{从“孙权亲攻合肥新城，魏守军与援军固守，吴军撤退”看见行动的边界}
\label{reading:tk-0234-wu-hefei-assault}

嘉禾三年（234年），账本所记的\eventanchor{TK-0234-WU-HEFEI-ASSAULT}{孙权亲攻合肥新城，魏守军与援军固守，吴军撤退}，首先不是抽象地图上移动的一枚棋子。读者应先把目光放在建业、武昌、濡须与长江港口：在这里，一道命令若要抵达，就必须穿过长江上下游、淮南接口、荆州峡口和长沙文书所见的地方。孙权亲攻合肥新城，魏守军与援军固守，吴军撤退。这一结果给出可核验的时间骨架；它同时提出更难的问题：谁能调动物资与人手，谁只能等待别人的决定，谁又在道路、水面或官署的缝隙里重新计算自己的位置。

从结构上说，舟师、港口、沿江城戍、州郡簿籍和能够随水路移动的粮食并不是储存在国库里便可自动使用的力量。它们需要在不同地点被登记、运出、交接、守护和补充；其中任何一环延误，远方的军队或使节就会面对与中枢不同的现实。蜀汉北伐迫使魏分兵关中，孙吴希望在淮南夺取可渡淮的据点并形成协同压力。这段前因不应被理解为单一人物预先设计的剧本，而是一组已经存在的条件：空间给出距离，制度分配发令权，地方社会承担把命令做成行动的成本。

行动者面对的选择因而很具体。对于这一政权或集团而言，关键不是笼统地“进”或“退”，而是把下游中枢、上游防线与继承安排维持在同一条可运转的江面上。若冒险推进，可能赢得一处节点、一个名义或一次短期主动；若保留力量，又可能把接口让给竞争者，或让既有同盟与部属怀疑中枢的能力。水军维持、岸防、疾病、远征补给和宫廷清洗互相放大不是事后的修辞，它在当时已构成选择的价格，只是材料很少留下所有人的意见与犹疑。

在这一限制下，孙权亲攻合肥新城，魏守军与援军固守，吴军撤退。把事实写到这里仍须克制：正史可以确认的通常是任命、出军、围守、退还、死亡、册封或政权转换等骨架；营中每一段传话、每一份账目和每一处战阵，往往没有同样坚实的记录。真正有解释力的，并非补足一个好看的传奇，而是观察这一骨架怎样被长江上下游、淮南接口、荆州峡口和长沙文书所见的地方的摩擦推向已经发生的结果。

直接效果首先落在可见的节点上。魏守住合肥防线，吴蜀并攻未能改变三国边界；两线同时作战的协调程度不可由后世战略推演假定。一个城、一个渡口、一个官署、一名辅政者或一条通道的状态改变，会立刻重排相邻行动者的风险：前线可能得到喘息，近畿可能获得声望，地方官却要重新面对征发、守备或接收的命令。事件的“结束”因此只是中央记事的一条句号；在运输、任官和地方关系中，它刚刚开始产生后果。

中期影响并不等于结果的线性放大。成功会使中枢相信某条路线、某位将领或某种制度安排值得继续投资，也可能把资源更加集中到已经拥挤的方向；失利则可能迫使它缩短战线、改写同盟或重排辅政权力。就此而言，禁军、辅政、太子网络、州郡与水陆指挥的连接既是此次行动的条件，也是它留下的制度遗产。国家能力并不是“有”或“无”，而是在每一次调用之后增添某种可重复的程序，同时制造另一处更难处理的依赖。

把镜头拉近，船只、粮米、劳役和地方登记使江面上的战略最终落到家庭与县廷。史书多从朝廷、将帅和使节写起，难以连续记录沿线家庭如何应付粮食、役使、迁徙和不确定的安全；这并不意味着他们只是背景。恰恰是这些无法被完整统计的劳动，使一条道路、一支舟师或一座城戍能够在若干季节内维持。读者不宜凭缺失的户口表虚构精确的数字，却应把这种材料缺口本身看作国家抵达社会并不均匀的证据。

这一事件也改变了与他者相处的方式。把下游中枢、上游防线与继承安排维持在同一条可运转的江面上在边境上经常意味着对方也必须改算自己的粮道、使节、驻军与继承风险。即使双方暂时不战，册命、招纳、投降或撤军也会传递有关实力与脆弱性的信息。外交在这里不是战争之外的礼节，而是以较低成本测试关系的一种行动；战争则是在谈判、承诺与防守网被证明不足之后，迫使对手承担更高成本的手段。

制度所得必须和制度损失一起衡量。此次行动若使命令链更清楚、边防节点更稳固、任官更能配合资源，它就扩大了中枢的可用能力；但若它迫使权力过度依赖少数近侍、边将、转输官或地方合作者，也会让下一次继承、叛乱或远征更难收束。水军维持、岸防、疾病、远征补给和宫廷清洗互相放大由此不只是军费或伤亡的代名词，更是信息集中、地方弹性与政治信任被逐步消耗的过程。

材料在这里决定叙述能走多远。本段的基本年序和行动依据是《三国志·吴主传》；《三国志·满宠传》；《三国志·明帝纪》；《资治通鉴》魏纪。它们能够相互校验一些结果，却不是互不相干的现场档案：纪传有编纂目的，后出的通鉴重组了时间，地方文献与考古又只照亮有限地点。因此，《三国志》吴书保存宫廷主线，走马楼吴简和区域考古仅能从局部校正国家叙事。引用多部书不自动制造全知视角；不同材料的沉默、重叠与矛盾，正是解释的一部分。

争议尤其集中在动机、数字与细节。地望、兵额、海上航程、宫廷对话和江东全境的户口都不能由一项材料裁定。较稳妥的写法是把可确认的结果与解释性的推断分开：可以说行动者在既有资源和压力下作出某种选择，也必须承认我们通常无法知道每个人在当时究竟相信什么、害怕什么。这样的保留并不会削弱历史，反而使读者看见选择为何有限、失败为何未必等于无能、胜利为何不必然证明先见。

如果把这一节点接回全章，它所显示的是一条反馈回路：舟师、港口、沿江城戍、州郡簿籍和能够随水路移动的粮食使政权能够作出行动；行动又把压力送回禁军、辅政、太子网络、州郡与水陆指挥的连接和船只、粮米、劳役和地方登记使江面上的战略最终落到家庭与县廷；新形成的权力关系反过来决定下一次资源由谁调用、沿哪条路线流动。三国史最值得阅读之处，正在于这种循环不断把宏大的名号压回具体的河流、山口、港口、文书与人事之中。孙权亲攻合肥新城，魏守军与援军固守，吴军撤退不是孤立条目，而是这条回路在嘉禾三年（234年）留下的一处可见折痕。


\subsection{胜负之外：禁军、辅政、太子网络、州郡与水陆指挥的连接的压力}
\label{reading:tk-0241-sun-deng-death}

赤乌四年（241年），账本所记的\eventanchor{TK-0241-SUN-DENG-DEATH}{太子孙登去世，孙吴失去已被长期培养的成年继承人}，首先不是抽象地图上移动的一枚棋子。读者应先把目光放在建业、武昌、濡须与长江港口：在这里，一道命令若要抵达，就必须穿过长江上下游、淮南接口、荆州峡口和长沙文书所见的地方。太子孙登去世，孙吴失去已被长期培养的成年继承人。这一结果给出可核验的时间骨架；它同时提出更难的问题：谁能调动物资与人手，谁只能等待别人的决定，谁又在道路、水面或官署的缝隙里重新计算自己的位置。

从结构上说，舟师、港口、沿江城戍、州郡簿籍和能够随水路移动的粮食并不是储存在国库里便可自动使用的力量。它们需要在不同地点被登记、运出、交接、守护和补充；其中任何一环延误，远方的军队或使节就会面对与中枢不同的现实。孙登原可在孙权晚年分担政务并平衡江东官僚、宗室与将领，其死亡使继承安排重新开放。这段前因不应被理解为单一人物预先设计的剧本，而是一组已经存在的条件：空间给出距离，制度分配发令权，地方社会承担把命令做成行动的成本。

行动者面对的选择因而很具体。对于这一政权或集团而言，关键不是笼统地“进”或“退”，而是把下游中枢、上游防线与继承安排维持在同一条可运转的江面上。若冒险推进，可能赢得一处节点、一个名义或一次短期主动；若保留力量，又可能把接口让给竞争者，或让既有同盟与部属怀疑中枢的能力。水军维持、岸防、疾病、远征补给和宫廷清洗互相放大不是事后的修辞，它在当时已构成选择的价格，只是材料很少留下所有人的意见与犹疑。

在这一限制下，太子孙登去世，孙吴失去已被长期培养的成年继承人。把事实写到这里仍须克制：正史可以确认的通常是任命、出军、围守、退还、死亡、册封或政权转换等骨架；营中每一段传话、每一份账目和每一处战阵，往往没有同样坚实的记录。真正有解释力的，并非补足一个好看的传奇，而是观察这一骨架怎样被长江上下游、淮南接口、荆州峡口和长沙文书所见的地方的摩擦推向已经发生的结果。

直接效果首先落在可见的节点上。孙权后来立孙和、封孙霸，储位竞争的条件形成；继承危机将削弱孙吴在淮南与荆州的决策稳定性。一个城、一个渡口、一个官署、一名辅政者或一条通道的状态改变，会立刻重排相邻行动者的风险：前线可能得到喘息，近畿可能获得声望，地方官却要重新面对征发、守备或接收的命令。事件的“结束”因此只是中央记事的一条句号；在运输、任官和地方关系中，它刚刚开始产生后果。

中期影响并不等于结果的线性放大。成功会使中枢相信某条路线、某位将领或某种制度安排值得继续投资，也可能把资源更加集中到已经拥挤的方向；失利则可能迫使它缩短战线、改写同盟或重排辅政权力。就此而言，禁军、辅政、太子网络、州郡与水陆指挥的连接既是此次行动的条件，也是它留下的制度遗产。国家能力并不是“有”或“无”，而是在每一次调用之后增添某种可重复的程序，同时制造另一处更难处理的依赖。

把镜头拉近，船只、粮米、劳役和地方登记使江面上的战略最终落到家庭与县廷。史书多从朝廷、将帅和使节写起，难以连续记录沿线家庭如何应付粮食、役使、迁徙和不确定的安全；这并不意味着他们只是背景。恰恰是这些无法被完整统计的劳动，使一条道路、一支舟师或一座城戍能够在若干季节内维持。读者不宜凭缺失的户口表虚构精确的数字，却应把这种材料缺口本身看作国家抵达社会并不均匀的证据。

这一事件也改变了与他者相处的方式。把下游中枢、上游防线与继承安排维持在同一条可运转的江面上在边境上经常意味着对方也必须改算自己的粮道、使节、驻军与继承风险。即使双方暂时不战，册命、招纳、投降或撤军也会传递有关实力与脆弱性的信息。外交在这里不是战争之外的礼节，而是以较低成本测试关系的一种行动；战争则是在谈判、承诺与防守网被证明不足之后，迫使对手承担更高成本的手段。

制度所得必须和制度损失一起衡量。此次行动若使命令链更清楚、边防节点更稳固、任官更能配合资源，它就扩大了中枢的可用能力；但若它迫使权力过度依赖少数近侍、边将、转输官或地方合作者，也会让下一次继承、叛乱或远征更难收束。水军维持、岸防、疾病、远征补给和宫廷清洗互相放大由此不只是军费或伤亡的代名词，更是信息集中、地方弹性与政治信任被逐步消耗的过程。

材料在这里决定叙述能走多远。本段的基本年序和行动依据是《三国志·孙登传》；《三国志·吴主传》；《建康实录》。它们能够相互校验一些结果，却不是互不相干的现场档案：纪传有编纂目的，后出的通鉴重组了时间，地方文献与考古又只照亮有限地点。因此，《三国志》吴书保存宫廷主线，走马楼吴简和区域考古仅能从局部校正国家叙事。引用多部书不自动制造全知视角；不同材料的沉默、重叠与矛盾，正是解释的一部分。

争议尤其集中在动机、数字与细节。地望、兵额、海上航程、宫廷对话和江东全境的户口都不能由一项材料裁定。较稳妥的写法是把可确认的结果与解释性的推断分开：可以说行动者在既有资源和压力下作出某种选择，也必须承认我们通常无法知道每个人在当时究竟相信什么、害怕什么。这样的保留并不会削弱历史，反而使读者看见选择为何有限、失败为何未必等于无能、胜利为何不必然证明先见。

如果把这一节点接回全章，它所显示的是一条反馈回路：舟师、港口、沿江城戍、州郡簿籍和能够随水路移动的粮食使政权能够作出行动；行动又把压力送回禁军、辅政、太子网络、州郡与水陆指挥的连接和船只、粮米、劳役和地方登记使江面上的战略最终落到家庭与县廷；新形成的权力关系反过来决定下一次资源由谁调用、沿哪条路线流动。三国史最值得阅读之处，正在于这种循环不断把宏大的名号压回具体的河流、山口、港口、文书与人事之中。太子孙登去世，孙吴失去已被长期培养的成年继承人不是孤立条目，而是这条回路在赤乌四年（241年）留下的一处可见折痕。


\subsection{建业、武昌、濡须与长江港口：一次转折的可达与不可达}
\label{reading:tk-0242-sun-he-sun-ba}

赤乌五年（242年），账本所记的\eventanchor{TK-0242-SUN-HE-SUN-BA}{孙和被立为太子，孙霸封鲁王，两个成年皇子各自形成宾客与官……}，首先不是抽象地图上移动的一枚棋子。读者应先把目光放在建业、武昌、濡须与长江港口：在这里，一道命令若要抵达，就必须穿过长江上下游、淮南接口、荆州峡口和长沙文书所见的地方。孙和被立为太子，孙霸封鲁王，两个成年皇子各自形成宾客与官僚支持网络。这一结果给出可核验的时间骨架；它同时提出更难的问题：谁能调动物资与人手，谁只能等待别人的决定，谁又在道路、水面或官署的缝隙里重新计算自己的位置。

从结构上说，舟师、港口、沿江城戍、州郡簿籍和能够随水路移动的粮食并不是储存在国库里便可自动使用的力量。它们需要在不同地点被登记、运出、交接、守护和补充；其中任何一环延误，远方的军队或使节就会面对与中枢不同的现实。孙登死亡后孙权需尽快确立继承人，却同时给予孙霸高位和可观的政治资源。这段前因不应被理解为单一人物预先设计的剧本，而是一组已经存在的条件：空间给出距离，制度分配发令权，地方社会承担把命令做成行动的成本。

行动者面对的选择因而很具体。对于这一政权或集团而言，关键不是笼统地“进”或“退”，而是把下游中枢、上游防线与继承安排维持在同一条可运转的江面上。若冒险推进，可能赢得一处节点、一个名义或一次短期主动；若保留力量，又可能把接口让给竞争者，或让既有同盟与部属怀疑中枢的能力。水军维持、岸防、疾病、远征补给和宫廷清洗互相放大不是事后的修辞，它在当时已构成选择的价格，只是材料很少留下所有人的意见与犹疑。

在这一限制下，孙和被立为太子，孙霸封鲁王，两个成年皇子各自形成宾客与官僚支持网络。把事实写到这里仍须克制：正史可以确认的通常是任命、出军、围守、退还、死亡、册封或政权转换等骨架；营中每一段传话、每一份账目和每一处战阵，往往没有同样坚实的记录。真正有解释力的，并非补足一个好看的传奇，而是观察这一骨架怎样被长江上下游、淮南接口、荆州峡口和长沙文书所见的地方的摩擦推向已经发生的结果。

直接效果首先落在可见的节点上。太子与鲁王并立使中枢官员被迫选边，孙权晚年的人事与军事决策将受到持续的储位冲突牵制。一个城、一个渡口、一个官署、一名辅政者或一条通道的状态改变，会立刻重排相邻行动者的风险：前线可能得到喘息，近畿可能获得声望，地方官却要重新面对征发、守备或接收的命令。事件的“结束”因此只是中央记事的一条句号；在运输、任官和地方关系中，它刚刚开始产生后果。

中期影响并不等于结果的线性放大。成功会使中枢相信某条路线、某位将领或某种制度安排值得继续投资，也可能把资源更加集中到已经拥挤的方向；失利则可能迫使它缩短战线、改写同盟或重排辅政权力。就此而言，禁军、辅政、太子网络、州郡与水陆指挥的连接既是此次行动的条件，也是它留下的制度遗产。国家能力并不是“有”或“无”，而是在每一次调用之后增添某种可重复的程序，同时制造另一处更难处理的依赖。

把镜头拉近，船只、粮米、劳役和地方登记使江面上的战略最终落到家庭与县廷。史书多从朝廷、将帅和使节写起，难以连续记录沿线家庭如何应付粮食、役使、迁徙和不确定的安全；这并不意味着他们只是背景。恰恰是这些无法被完整统计的劳动，使一条道路、一支舟师或一座城戍能够在若干季节内维持。读者不宜凭缺失的户口表虚构精确的数字，却应把这种材料缺口本身看作国家抵达社会并不均匀的证据。

这一事件也改变了与他者相处的方式。把下游中枢、上游防线与继承安排维持在同一条可运转的江面上在边境上经常意味着对方也必须改算自己的粮道、使节、驻军与继承风险。即使双方暂时不战，册命、招纳、投降或撤军也会传递有关实力与脆弱性的信息。外交在这里不是战争之外的礼节，而是以较低成本测试关系的一种行动；战争则是在谈判、承诺与防守网被证明不足之后，迫使对手承担更高成本的手段。

制度所得必须和制度损失一起衡量。此次行动若使命令链更清楚、边防节点更稳固、任官更能配合资源，它就扩大了中枢的可用能力；但若它迫使权力过度依赖少数近侍、边将、转输官或地方合作者，也会让下一次继承、叛乱或远征更难收束。水军维持、岸防、疾病、远征补给和宫廷清洗互相放大由此不只是军费或伤亡的代名词，更是信息集中、地方弹性与政治信任被逐步消耗的过程。

材料在这里决定叙述能走多远。本段的基本年序和行动依据是《三国志·孙和传》；《三国志·孙霸传》裴松之注引《吴录》；《三国志·吴主传》。它们能够相互校验一些结果，却不是互不相干的现场档案：纪传有编纂目的，后出的通鉴重组了时间，地方文献与考古又只照亮有限地点。因此，《三国志》吴书保存宫廷主线，走马楼吴简和区域考古仅能从局部校正国家叙事。引用多部书不自动制造全知视角；不同材料的沉默、重叠与矛盾，正是解释的一部分。

争议尤其集中在动机、数字与细节。地望、兵额、海上航程、宫廷对话和江东全境的户口都不能由一项材料裁定。较稳妥的写法是把可确认的结果与解释性的推断分开：可以说行动者在既有资源和压力下作出某种选择，也必须承认我们通常无法知道每个人在当时究竟相信什么、害怕什么。这样的保留并不会削弱历史，反而使读者看见选择为何有限、失败为何未必等于无能、胜利为何不必然证明先见。

如果把这一节点接回全章，它所显示的是一条反馈回路：舟师、港口、沿江城戍、州郡簿籍和能够随水路移动的粮食使政权能够作出行动；行动又把压力送回禁军、辅政、太子网络、州郡与水陆指挥的连接和船只、粮米、劳役和地方登记使江面上的战略最终落到家庭与县廷；新形成的权力关系反过来决定下一次资源由谁调用、沿哪条路线流动。三国史最值得阅读之处，正在于这种循环不断把宏大的名号压回具体的河流、山口、港口、文书与人事之中。孙和被立为太子，孙霸封鲁王，两个成年皇子各自形成宾客与官僚支持网络不是孤立条目，而是这条回路在赤乌五年（242年）留下的一处可见折痕。

\begin{sourcenote}
这一组事件的来源路径并不相同。《三国志》吴书保存宫廷主线，走马楼吴简和区域考古仅能从局部校正国家叙事。阅读时应先区分能确认的年序和结果，与只能作为解释线索的动机、数量及局部过程；这比用一个完整无缝的故事覆盖材料间隙更能说明国家能力如何实际运作。
\end{sourcenote}


\subsection{建业、武昌、濡须与长江港口：一次转折的可达与不可达}
\label{reading:tk-0245-lu-xun-death}

赤乌八年（245年），账本所记的\eventanchor{TK-0245-LU-XUN-DEATH}{丞相陆逊去世，其与孙权围绕太子孙和、鲁王孙霸的谏争成为后……}，首先不是抽象地图上移动的一枚棋子。读者应先把目光放在建业、武昌、濡须与长江港口：在这里，一道命令若要抵达，就必须穿过长江上下游、淮南接口、荆州峡口和长沙文书所见的地方。丞相陆逊去世，其与孙权围绕太子孙和、鲁王孙霸的谏争成为后世叙述焦点。这一结果给出可核验的时间骨架；它同时提出更难的问题：谁能调动物资与人手，谁只能等待别人的决定，谁又在道路、水面或官署的缝隙里重新计算自己的位置。

从结构上说，舟师、港口、沿江城戍、州郡簿籍和能够随水路移动的粮食并不是储存在国库里便可自动使用的力量。它们需要在不同地点被登记、运出、交接、守护和补充；其中任何一环延误，远方的军队或使节就会面对与中枢不同的现实。二宫党争侵入军政人事，陆逊以丞相和武昌方面重臣身份干预储位，触及孙权晚年的决断。这段前因不应被理解为单一人物预先设计的剧本，而是一组已经存在的条件：空间给出距离，制度分配发令权，地方社会承担把命令做成行动的成本。

行动者面对的选择因而很具体。对于这一政权或集团而言，关键不是笼统地“进”或“退”，而是把下游中枢、上游防线与继承安排维持在同一条可运转的江面上。若冒险推进，可能赢得一处节点、一个名义或一次短期主动；若保留力量，又可能把接口让给竞争者，或让既有同盟与部属怀疑中枢的能力。水军维持、岸防、疾病、远征补给和宫廷清洗互相放大不是事后的修辞，它在当时已构成选择的价格，只是材料很少留下所有人的意见与犹疑。

在这一限制下，丞相陆逊去世，其与孙权围绕太子孙和、鲁王孙霸的谏争成为后世叙述焦点。把事实写到这里仍须克制：正史可以确认的通常是任命、出军、围守、退还、死亡、册封或政权转换等骨架；营中每一段传话、每一份账目和每一处战阵，往往没有同样坚实的记录。真正有解释力的，并非补足一个好看的传奇，而是观察这一骨架怎样被长江上下游、淮南接口、荆州峡口和长沙文书所见的地方的摩擦推向已经发生的结果。

直接效果首先落在可见的节点上。孙吴失去兼具荆州军政经验的重臣，储位冲突更难调停；后出书信需要辨析其传抄和政治立场。一个城、一个渡口、一个官署、一名辅政者或一条通道的状态改变，会立刻重排相邻行动者的风险：前线可能得到喘息，近畿可能获得声望，地方官却要重新面对征发、守备或接收的命令。事件的“结束”因此只是中央记事的一条句号；在运输、任官和地方关系中，它刚刚开始产生后果。

中期影响并不等于结果的线性放大。成功会使中枢相信某条路线、某位将领或某种制度安排值得继续投资，也可能把资源更加集中到已经拥挤的方向；失利则可能迫使它缩短战线、改写同盟或重排辅政权力。就此而言，禁军、辅政、太子网络、州郡与水陆指挥的连接既是此次行动的条件，也是它留下的制度遗产。国家能力并不是“有”或“无”，而是在每一次调用之后增添某种可重复的程序，同时制造另一处更难处理的依赖。

把镜头拉近，船只、粮米、劳役和地方登记使江面上的战略最终落到家庭与县廷。史书多从朝廷、将帅和使节写起，难以连续记录沿线家庭如何应付粮食、役使、迁徙和不确定的安全；这并不意味着他们只是背景。恰恰是这些无法被完整统计的劳动，使一条道路、一支舟师或一座城戍能够在若干季节内维持。读者不宜凭缺失的户口表虚构精确的数字，却应把这种材料缺口本身看作国家抵达社会并不均匀的证据。

这一事件也改变了与他者相处的方式。把下游中枢、上游防线与继承安排维持在同一条可运转的江面上在边境上经常意味着对方也必须改算自己的粮道、使节、驻军与继承风险。即使双方暂时不战，册命、招纳、投降或撤军也会传递有关实力与脆弱性的信息。外交在这里不是战争之外的礼节，而是以较低成本测试关系的一种行动；战争则是在谈判、承诺与防守网被证明不足之后，迫使对手承担更高成本的手段。

制度所得必须和制度损失一起衡量。此次行动若使命令链更清楚、边防节点更稳固、任官更能配合资源，它就扩大了中枢的可用能力；但若它迫使权力过度依赖少数近侍、边将、转输官或地方合作者，也会让下一次继承、叛乱或远征更难收束。水军维持、岸防、疾病、远征补给和宫廷清洗互相放大由此不只是军费或伤亡的代名词，更是信息集中、地方弹性与政治信任被逐步消耗的过程。

材料在这里决定叙述能走多远。本段的基本年序和行动依据是《三国志·陆逊传》；《三国志·孙和传》；《三国志·陆逊传》裴松之注引《吴书》。它们能够相互校验一些结果，却不是互不相干的现场档案：纪传有编纂目的，后出的通鉴重组了时间，地方文献与考古又只照亮有限地点。因此，《三国志》吴书保存宫廷主线，走马楼吴简和区域考古仅能从局部校正国家叙事。引用多部书不自动制造全知视角；不同材料的沉默、重叠与矛盾，正是解释的一部分。

争议尤其集中在动机、数字与细节。地望、兵额、海上航程、宫廷对话和江东全境的户口都不能由一项材料裁定。较稳妥的写法是把可确认的结果与解释性的推断分开：可以说行动者在既有资源和压力下作出某种选择，也必须承认我们通常无法知道每个人在当时究竟相信什么、害怕什么。这样的保留并不会削弱历史，反而使读者看见选择为何有限、失败为何未必等于无能、胜利为何不必然证明先见。

如果把这一节点接回全章，它所显示的是一条反馈回路：舟师、港口、沿江城戍、州郡簿籍和能够随水路移动的粮食使政权能够作出行动；行动又把压力送回禁军、辅政、太子网络、州郡与水陆指挥的连接和船只、粮米、劳役和地方登记使江面上的战略最终落到家庭与县廷；新形成的权力关系反过来决定下一次资源由谁调用、沿哪条路线流动。三国史最值得阅读之处，正在于这种循环不断把宏大的名号压回具体的河流、山口、港口、文书与人事之中。丞相陆逊去世，其与孙权围绕太子孙和、鲁王孙霸的谏争成为后世叙述焦点不是孤立条目，而是这条回路在赤乌八年（245年）留下的一处可见折痕。


\subsection{建业、武昌、濡须与长江港口：一次转折的可达与不可达}
\label{reading:tk-0250-sun-wu-succession-crisis}

赤乌十三年（250年），账本所记的\eventanchor{TK-0250-SUN-WU-SUCCESSION-CRISIS}{孙权废太子孙和、赐死鲁王孙霸，改立幼子孙亮为太子，二宫之……}，首先不是抽象地图上移动的一枚棋子。读者应先把目光放在建业、武昌、濡须与长江港口：在这里，一道命令若要抵达，就必须穿过长江上下游、淮南接口、荆州峡口和长沙文书所见的地方。孙权废太子孙和、赐死鲁王孙霸，改立幼子孙亮为太子，二宫之争以强制清洗收束。这一结果给出可核验的时间骨架；它同时提出更难的问题：谁能调动物资与人手，谁只能等待别人的决定，谁又在道路、水面或官署的缝隙里重新计算自己的位置。

从结构上说，舟师、港口、沿江城戍、州郡簿籍和能够随水路移动的粮食并不是储存在国库里便可自动使用的力量。它们需要在不同地点被登记、运出、交接、守护和补充；其中任何一环延误，远方的军队或使节就会面对与中枢不同的现实。孙和与孙霸各自形成的宾客集团长期侵蚀朝廷，孙权既不能容忍持续分裂，也缺少能兼得各方认可的成年继承人。这段前因不应被理解为单一人物预先设计的剧本，而是一组已经存在的条件：空间给出距离，制度分配发令权，地方社会承担把命令做成行动的成本。

行动者面对的选择因而很具体。对于这一政权或集团而言，关键不是笼统地“进”或“退”，而是把下游中枢、上游防线与继承安排维持在同一条可运转的江面上。若冒险推进，可能赢得一处节点、一个名义或一次短期主动；若保留力量，又可能把接口让给竞争者，或让既有同盟与部属怀疑中枢的能力。水军维持、岸防、疾病、远征补给和宫廷清洗互相放大不是事后的修辞，它在当时已构成选择的价格，只是材料很少留下所有人的意见与犹疑。

在这一限制下，孙权废太子孙和、赐死鲁王孙霸，改立幼子孙亮为太子，二宫之争以强制清洗收束。把事实写到这里仍须克制：正史可以确认的通常是任命、出军、围守、退还、死亡、册封或政权转换等骨架；营中每一段传话、每一份账目和每一处战阵，往往没有同样坚实的记录。真正有解释力的，并非补足一个好看的传奇，而是观察这一骨架怎样被长江上下游、淮南接口、荆州峡口和长沙文书所见的地方的摩擦推向已经发生的结果。

直接效果首先落在可见的节点上。大量官员被诛逐，孙亮继承使孙吴再次面对幼主与权臣问题；事件不能被化约为单纯的家族伦理冲突。一个城、一个渡口、一个官署、一名辅政者或一条通道的状态改变，会立刻重排相邻行动者的风险：前线可能得到喘息，近畿可能获得声望，地方官却要重新面对征发、守备或接收的命令。事件的“结束”因此只是中央记事的一条句号；在运输、任官和地方关系中，它刚刚开始产生后果。

中期影响并不等于结果的线性放大。成功会使中枢相信某条路线、某位将领或某种制度安排值得继续投资，也可能把资源更加集中到已经拥挤的方向；失利则可能迫使它缩短战线、改写同盟或重排辅政权力。就此而言，禁军、辅政、太子网络、州郡与水陆指挥的连接既是此次行动的条件，也是它留下的制度遗产。国家能力并不是“有”或“无”，而是在每一次调用之后增添某种可重复的程序，同时制造另一处更难处理的依赖。

把镜头拉近，船只、粮米、劳役和地方登记使江面上的战略最终落到家庭与县廷。史书多从朝廷、将帅和使节写起，难以连续记录沿线家庭如何应付粮食、役使、迁徙和不确定的安全；这并不意味着他们只是背景。恰恰是这些无法被完整统计的劳动，使一条道路、一支舟师或一座城戍能够在若干季节内维持。读者不宜凭缺失的户口表虚构精确的数字，却应把这种材料缺口本身看作国家抵达社会并不均匀的证据。

这一事件也改变了与他者相处的方式。把下游中枢、上游防线与继承安排维持在同一条可运转的江面上在边境上经常意味着对方也必须改算自己的粮道、使节、驻军与继承风险。即使双方暂时不战，册命、招纳、投降或撤军也会传递有关实力与脆弱性的信息。外交在这里不是战争之外的礼节，而是以较低成本测试关系的一种行动；战争则是在谈判、承诺与防守网被证明不足之后，迫使对手承担更高成本的手段。

制度所得必须和制度损失一起衡量。此次行动若使命令链更清楚、边防节点更稳固、任官更能配合资源，它就扩大了中枢的可用能力；但若它迫使权力过度依赖少数近侍、边将、转输官或地方合作者，也会让下一次继承、叛乱或远征更难收束。水军维持、岸防、疾病、远征补给和宫廷清洗互相放大由此不只是军费或伤亡的代名词，更是信息集中、地方弹性与政治信任被逐步消耗的过程。

材料在这里决定叙述能走多远。本段的基本年序和行动依据是《三国志·孙和传》；《三国志·孙亮传》；《三国志·吴主传》；《资治通鉴》魏纪。它们能够相互校验一些结果，却不是互不相干的现场档案：纪传有编纂目的，后出的通鉴重组了时间，地方文献与考古又只照亮有限地点。因此，《三国志》吴书保存宫廷主线，走马楼吴简和区域考古仅能从局部校正国家叙事。引用多部书不自动制造全知视角；不同材料的沉默、重叠与矛盾，正是解释的一部分。

争议尤其集中在动机、数字与细节。地望、兵额、海上航程、宫廷对话和江东全境的户口都不能由一项材料裁定。较稳妥的写法是把可确认的结果与解释性的推断分开：可以说行动者在既有资源和压力下作出某种选择，也必须承认我们通常无法知道每个人在当时究竟相信什么、害怕什么。这样的保留并不会削弱历史，反而使读者看见选择为何有限、失败为何未必等于无能、胜利为何不必然证明先见。

如果把这一节点接回全章，它所显示的是一条反馈回路：舟师、港口、沿江城戍、州郡簿籍和能够随水路移动的粮食使政权能够作出行动；行动又把压力送回禁军、辅政、太子网络、州郡与水陆指挥的连接和船只、粮米、劳役和地方登记使江面上的战略最终落到家庭与县廷；新形成的权力关系反过来决定下一次资源由谁调用、沿哪条路线流动。三国史最值得阅读之处，正在于这种循环不断把宏大的名号压回具体的河流、山口、港口、文书与人事之中。孙权废太子孙和、赐死鲁王孙霸，改立幼子孙亮为太子，二宫之争以强制清洗收束不是孤立条目，而是这条回路在赤乌十三年（250年）留下的一处可见折痕。


\subsection{胜负之外：禁军、辅政、太子网络、州郡与水陆指挥的连接的压力}
\label{reading:tk-0252-sun-quan-death}

太元二年（252年），账本所记的\eventanchor{TK-0252-SUN-QUAN-DEATH}{孙权去世，太子孙亮继位，诸葛恪受遗诏辅政}，首先不是抽象地图上移动的一枚棋子。读者应先把目光放在建业、武昌、濡须与长江港口：在这里，一道命令若要抵达，就必须穿过长江上下游、淮南接口、荆州峡口和长沙文书所见的地方。孙权去世，太子孙亮继位，诸葛恪受遗诏辅政。这一结果给出可核验的时间骨架；它同时提出更难的问题：谁能调动物资与人手，谁只能等待别人的决定，谁又在道路、水面或官署的缝隙里重新计算自己的位置。

从结构上说，舟师、港口、沿江城戍、州郡簿籍和能够随水路移动的粮食并不是储存在国库里便可自动使用的力量。它们需要在不同地点被登记、运出、交接、守护和补充；其中任何一环延误，远方的军队或使节就会面对与中枢不同的现实。二宫之争已清洗大批重臣，孙亮年幼，孙权只能依靠外戚性重臣和宗室军人维持中枢与边防。这段前因不应被理解为单一人物预先设计的剧本，而是一组已经存在的条件：空间给出距离，制度分配发令权，地方社会承担把命令做成行动的成本。

行动者面对的选择因而很具体。对于这一政权或集团而言，关键不是笼统地“进”或“退”，而是把下游中枢、上游防线与继承安排维持在同一条可运转的江面上。若冒险推进，可能赢得一处节点、一个名义或一次短期主动；若保留力量，又可能把接口让给竞争者，或让既有同盟与部属怀疑中枢的能力。水军维持、岸防、疾病、远征补给和宫廷清洗互相放大不是事后的修辞，它在当时已构成选择的价格，只是材料很少留下所有人的意见与犹疑。

在这一限制下，孙权去世，太子孙亮继位，诸葛恪受遗诏辅政。把事实写到这里仍须克制：正史可以确认的通常是任命、出军、围守、退还、死亡、册封或政权转换等骨架；营中每一段传话、每一份账目和每一处战阵，往往没有同样坚实的记录。真正有解释力的，并非补足一个好看的传奇，而是观察这一骨架怎样被长江上下游、淮南接口、荆州峡口和长沙文书所见的地方的摩擦推向已经发生的结果。

直接效果首先落在可见的节点上。孙吴进入诸葛恪主政期，辅政集团的军事胜负将直接影响幼帝和江东官僚对其合法性的判断。一个城、一个渡口、一个官署、一名辅政者或一条通道的状态改变，会立刻重排相邻行动者的风险：前线可能得到喘息，近畿可能获得声望，地方官却要重新面对征发、守备或接收的命令。事件的“结束”因此只是中央记事的一条句号；在运输、任官和地方关系中，它刚刚开始产生后果。

中期影响并不等于结果的线性放大。成功会使中枢相信某条路线、某位将领或某种制度安排值得继续投资，也可能把资源更加集中到已经拥挤的方向；失利则可能迫使它缩短战线、改写同盟或重排辅政权力。就此而言，禁军、辅政、太子网络、州郡与水陆指挥的连接既是此次行动的条件，也是它留下的制度遗产。国家能力并不是“有”或“无”，而是在每一次调用之后增添某种可重复的程序，同时制造另一处更难处理的依赖。

把镜头拉近，船只、粮米、劳役和地方登记使江面上的战略最终落到家庭与县廷。史书多从朝廷、将帅和使节写起，难以连续记录沿线家庭如何应付粮食、役使、迁徙和不确定的安全；这并不意味着他们只是背景。恰恰是这些无法被完整统计的劳动，使一条道路、一支舟师或一座城戍能够在若干季节内维持。读者不宜凭缺失的户口表虚构精确的数字，却应把这种材料缺口本身看作国家抵达社会并不均匀的证据。

这一事件也改变了与他者相处的方式。把下游中枢、上游防线与继承安排维持在同一条可运转的江面上在边境上经常意味着对方也必须改算自己的粮道、使节、驻军与继承风险。即使双方暂时不战，册命、招纳、投降或撤军也会传递有关实力与脆弱性的信息。外交在这里不是战争之外的礼节，而是以较低成本测试关系的一种行动；战争则是在谈判、承诺与防守网被证明不足之后，迫使对手承担更高成本的手段。

制度所得必须和制度损失一起衡量。此次行动若使命令链更清楚、边防节点更稳固、任官更能配合资源，它就扩大了中枢的可用能力；但若它迫使权力过度依赖少数近侍、边将、转输官或地方合作者，也会让下一次继承、叛乱或远征更难收束。水军维持、岸防、疾病、远征补给和宫廷清洗互相放大由此不只是军费或伤亡的代名词，更是信息集中、地方弹性与政治信任被逐步消耗的过程。

材料在这里决定叙述能走多远。本段的基本年序和行动依据是《三国志·吴主传》；《三国志·孙亮传》；《三国志·诸葛恪传》。它们能够相互校验一些结果，却不是互不相干的现场档案：纪传有编纂目的，后出的通鉴重组了时间，地方文献与考古又只照亮有限地点。因此，《三国志》吴书保存宫廷主线，走马楼吴简和区域考古仅能从局部校正国家叙事。引用多部书不自动制造全知视角；不同材料的沉默、重叠与矛盾，正是解释的一部分。

争议尤其集中在动机、数字与细节。地望、兵额、海上航程、宫廷对话和江东全境的户口都不能由一项材料裁定。较稳妥的写法是把可确认的结果与解释性的推断分开：可以说行动者在既有资源和压力下作出某种选择，也必须承认我们通常无法知道每个人在当时究竟相信什么、害怕什么。这样的保留并不会削弱历史，反而使读者看见选择为何有限、失败为何未必等于无能、胜利为何不必然证明先见。

如果把这一节点接回全章，它所显示的是一条反馈回路：舟师、港口、沿江城戍、州郡簿籍和能够随水路移动的粮食使政权能够作出行动；行动又把压力送回禁军、辅政、太子网络、州郡与水陆指挥的连接和船只、粮米、劳役和地方登记使江面上的战略最终落到家庭与县廷；新形成的权力关系反过来决定下一次资源由谁调用、沿哪条路线流动。三国史最值得阅读之处，正在于这种循环不断把宏大的名号压回具体的河流、山口、港口、文书与人事之中。孙权去世，太子孙亮继位，诸葛恪受遗诏辅政不是孤立条目，而是这条回路在太元二年（252年）留下的一处可见折痕。


\subsection{建业、武昌、濡须与长江港口：一次转折的可达与不可达}
\label{reading:tk-0252-dongxing}

建兴元年（252年），账本所记的\eventanchor{TK-0252-DONGXING}{诸葛恪在东兴堤击退魏军，吴军利用水陆地形保住濡须方向}，首先不是抽象地图上移动的一枚棋子。读者应先把目光放在建业、武昌、濡须与长江港口：在这里，一道命令若要抵达，就必须穿过长江上下游、淮南接口、荆州峡口和长沙文书所见的地方。诸葛恪在东兴堤击退魏军，吴军利用水陆地形保住濡须方向。这一结果给出可核验的时间骨架；它同时提出更难的问题：谁能调动物资与人手，谁只能等待别人的决定，谁又在道路、水面或官署的缝隙里重新计算自己的位置。

从结构上说，舟师、港口、沿江城戍、州郡簿籍和能够随水路移动的粮食并不是储存在国库里便可自动使用的力量。它们需要在不同地点被登记、运出、交接、守护和补充；其中任何一环延误，远方的军队或使节就会面对与中枢不同的现实。孙权死后魏试图以大军试探幼主政权，东兴堤和湖泊水网成为双方争夺的前沿据点。这段前因不应被理解为单一人物预先设计的剧本，而是一组已经存在的条件：空间给出距离，制度分配发令权，地方社会承担把命令做成行动的成本。

行动者面对的选择因而很具体。对于这一政权或集团而言，关键不是笼统地“进”或“退”，而是把下游中枢、上游防线与继承安排维持在同一条可运转的江面上。若冒险推进，可能赢得一处节点、一个名义或一次短期主动；若保留力量，又可能把接口让给竞争者，或让既有同盟与部属怀疑中枢的能力。水军维持、岸防、疾病、远征补给和宫廷清洗互相放大不是事后的修辞，它在当时已构成选择的价格，只是材料很少留下所有人的意见与犹疑。

在这一限制下，诸葛恪在东兴堤击退魏军，吴军利用水陆地形保住濡须方向。把事实写到这里仍须克制：正史可以确认的通常是任命、出军、围守、退还、死亡、册封或政权转换等骨架；营中每一段传话、每一份账目和每一处战阵，往往没有同样坚实的记录。真正有解释力的，并非补足一个好看的传奇，而是观察这一骨架怎样被长江上下游、淮南接口、荆州峡口和长沙文书所见的地方的摩擦推向已经发生的结果。

直接效果首先落在可见的节点上。孙吴取得诸葛恪辅政后的重要胜利，短期提振江东士气；传记的突击叙事不宜直接换算为精确战损。一个城、一个渡口、一个官署、一名辅政者或一条通道的状态改变，会立刻重排相邻行动者的风险：前线可能得到喘息，近畿可能获得声望，地方官却要重新面对征发、守备或接收的命令。事件的“结束”因此只是中央记事的一条句号；在运输、任官和地方关系中，它刚刚开始产生后果。

中期影响并不等于结果的线性放大。成功会使中枢相信某条路线、某位将领或某种制度安排值得继续投资，也可能把资源更加集中到已经拥挤的方向；失利则可能迫使它缩短战线、改写同盟或重排辅政权力。就此而言，禁军、辅政、太子网络、州郡与水陆指挥的连接既是此次行动的条件，也是它留下的制度遗产。国家能力并不是“有”或“无”，而是在每一次调用之后增添某种可重复的程序，同时制造另一处更难处理的依赖。

把镜头拉近，船只、粮米、劳役和地方登记使江面上的战略最终落到家庭与县廷。史书多从朝廷、将帅和使节写起，难以连续记录沿线家庭如何应付粮食、役使、迁徙和不确定的安全；这并不意味着他们只是背景。恰恰是这些无法被完整统计的劳动，使一条道路、一支舟师或一座城戍能够在若干季节内维持。读者不宜凭缺失的户口表虚构精确的数字，却应把这种材料缺口本身看作国家抵达社会并不均匀的证据。

这一事件也改变了与他者相处的方式。把下游中枢、上游防线与继承安排维持在同一条可运转的江面上在边境上经常意味着对方也必须改算自己的粮道、使节、驻军与继承风险。即使双方暂时不战，册命、招纳、投降或撤军也会传递有关实力与脆弱性的信息。外交在这里不是战争之外的礼节，而是以较低成本测试关系的一种行动；战争则是在谈判、承诺与防守网被证明不足之后，迫使对手承担更高成本的手段。

制度所得必须和制度损失一起衡量。此次行动若使命令链更清楚、边防节点更稳固、任官更能配合资源，它就扩大了中枢的可用能力；但若它迫使权力过度依赖少数近侍、边将、转输官或地方合作者，也会让下一次继承、叛乱或远征更难收束。水军维持、岸防、疾病、远征补给和宫廷清洗互相放大由此不只是军费或伤亡的代名词，更是信息集中、地方弹性与政治信任被逐步消耗的过程。

材料在这里决定叙述能走多远。本段的基本年序和行动依据是《三国志·诸葛恪传》；《三国志·丁奉传》；《三国志·齐王芳纪》；《资治通鉴》魏纪。它们能够相互校验一些结果，却不是互不相干的现场档案：纪传有编纂目的，后出的通鉴重组了时间，地方文献与考古又只照亮有限地点。因此，《三国志》吴书保存宫廷主线，走马楼吴简和区域考古仅能从局部校正国家叙事。引用多部书不自动制造全知视角；不同材料的沉默、重叠与矛盾，正是解释的一部分。

争议尤其集中在动机、数字与细节。地望、兵额、海上航程、宫廷对话和江东全境的户口都不能由一项材料裁定。较稳妥的写法是把可确认的结果与解释性的推断分开：可以说行动者在既有资源和压力下作出某种选择，也必须承认我们通常无法知道每个人在当时究竟相信什么、害怕什么。这样的保留并不会削弱历史，反而使读者看见选择为何有限、失败为何未必等于无能、胜利为何不必然证明先见。

如果把这一节点接回全章，它所显示的是一条反馈回路：舟师、港口、沿江城戍、州郡簿籍和能够随水路移动的粮食使政权能够作出行动；行动又把压力送回禁军、辅政、太子网络、州郡与水陆指挥的连接和船只、粮米、劳役和地方登记使江面上的战略最终落到家庭与县廷；新形成的权力关系反过来决定下一次资源由谁调用、沿哪条路线流动。三国史最值得阅读之处，正在于这种循环不断把宏大的名号压回具体的河流、山口、港口、文书与人事之中。诸葛恪在东兴堤击退魏军，吴军利用水陆地形保住濡须方向不是孤立条目，而是这条回路在建兴元年（252年）留下的一处可见折痕。

\begin{sourcenote}
这一组事件的来源路径并不相同。《三国志》吴书保存宫廷主线，走马楼吴简和区域考古仅能从局部校正国家叙事。阅读时应先区分能确认的年序和结果，与只能作为解释线索的动机、数量及局部过程；这比用一个完整无缝的故事覆盖材料间隙更能说明国家能力如何实际运作。
\end{sourcenote}


\subsection{胜负之外：禁军、辅政、太子网络、州郡与水陆指挥的连接的压力}
\label{reading:tk-0253-zhuge-ke-hefei}

建兴二年（253年），账本所记的\eventref{TK-0253-ZHUGE-KE-HEFEI}{诸葛恪围攻合肥新城久攻不下，吴军因疫病、补给与魏守备而撤退}，首先不是抽象地图上移动的一枚棋子。读者应先把目光放在建业、武昌、濡须与长江港口：在这里，一道命令若要抵达，就必须穿过长江上下游、淮南接口、荆州峡口和长沙文书所见的地方。诸葛恪围攻合肥新城久攻不下，吴军因疫病、补给与魏守备而撤退。这一结果给出可核验的时间骨架；它同时提出更难的问题：谁能调动物资与人手，谁只能等待别人的决定，谁又在道路、水面或官署的缝隙里重新计算自己的位置。

从结构上说，舟师、港口、沿江城戍、州郡簿籍和能够随水路移动的粮食并不是储存在国库里便可自动使用的力量。它们需要在不同地点被登记、运出、交接、守护和补充；其中任何一环延误，远方的军队或使节就会面对与中枢不同的现实。东兴胜利提高诸葛恪威望，他试图趁势夺取合肥以改变淮南战略格局，却低估了攻城补给与夏季疾病成本。这段前因不应被理解为单一人物预先设计的剧本，而是一组已经存在的条件：空间给出距离，制度分配发令权，地方社会承担把命令做成行动的成本。

行动者面对的选择因而很具体。对于这一政权或集团而言，关键不是笼统地“进”或“退”，而是把下游中枢、上游防线与继承安排维持在同一条可运转的江面上。若冒险推进，可能赢得一处节点、一个名义或一次短期主动；若保留力量，又可能把接口让给竞争者，或让既有同盟与部属怀疑中枢的能力。水军维持、岸防、疾病、远征补给和宫廷清洗互相放大不是事后的修辞，它在当时已构成选择的价格，只是材料很少留下所有人的意见与犹疑。

在这一限制下，诸葛恪围攻合肥新城久攻不下，吴军因疫病、补给与魏守备而撤退。把事实写到这里仍须克制：正史可以确认的通常是任命、出军、围守、退还、死亡、册封或政权转换等骨架；营中每一段传话、每一份账目和每一处战阵，往往没有同样坚实的记录。真正有解释力的，并非补足一个好看的传奇，而是观察这一骨架怎样被长江上下游、淮南接口、荆州峡口和长沙文书所见的地方的摩擦推向已经发生的结果。

直接效果首先落在可见的节点上。远征损害吴军和诸葛恪的政治信誉，为其被孙峻等人推翻创造条件；疾病规模不能由史书修辞量化。一个城、一个渡口、一个官署、一名辅政者或一条通道的状态改变，会立刻重排相邻行动者的风险：前线可能得到喘息，近畿可能获得声望，地方官却要重新面对征发、守备或接收的命令。事件的“结束”因此只是中央记事的一条句号；在运输、任官和地方关系中，它刚刚开始产生后果。

中期影响并不等于结果的线性放大。成功会使中枢相信某条路线、某位将领或某种制度安排值得继续投资，也可能把资源更加集中到已经拥挤的方向；失利则可能迫使它缩短战线、改写同盟或重排辅政权力。就此而言，禁军、辅政、太子网络、州郡与水陆指挥的连接既是此次行动的条件，也是它留下的制度遗产。国家能力并不是“有”或“无”，而是在每一次调用之后增添某种可重复的程序，同时制造另一处更难处理的依赖。

把镜头拉近，船只、粮米、劳役和地方登记使江面上的战略最终落到家庭与县廷。史书多从朝廷、将帅和使节写起，难以连续记录沿线家庭如何应付粮食、役使、迁徙和不确定的安全；这并不意味着他们只是背景。恰恰是这些无法被完整统计的劳动，使一条道路、一支舟师或一座城戍能够在若干季节内维持。读者不宜凭缺失的户口表虚构精确的数字，却应把这种材料缺口本身看作国家抵达社会并不均匀的证据。

这一事件也改变了与他者相处的方式。把下游中枢、上游防线与继承安排维持在同一条可运转的江面上在边境上经常意味着对方也必须改算自己的粮道、使节、驻军与继承风险。即使双方暂时不战，册命、招纳、投降或撤军也会传递有关实力与脆弱性的信息。外交在这里不是战争之外的礼节，而是以较低成本测试关系的一种行动；战争则是在谈判、承诺与防守网被证明不足之后，迫使对手承担更高成本的手段。

制度所得必须和制度损失一起衡量。此次行动若使命令链更清楚、边防节点更稳固、任官更能配合资源，它就扩大了中枢的可用能力；但若它迫使权力过度依赖少数近侍、边将、转输官或地方合作者，也会让下一次继承、叛乱或远征更难收束。水军维持、岸防、疾病、远征补给和宫廷清洗互相放大由此不只是军费或伤亡的代名词，更是信息集中、地方弹性与政治信任被逐步消耗的过程。

材料在这里决定叙述能走多远。本段的基本年序和行动依据是《三国志·诸葛恪传》；《三国志·孙亮传》；《三国志·毌丘俭传》；《资治通鉴》魏纪。它们能够相互校验一些结果，却不是互不相干的现场档案：纪传有编纂目的，后出的通鉴重组了时间，地方文献与考古又只照亮有限地点。因此，《三国志》吴书保存宫廷主线，走马楼吴简和区域考古仅能从局部校正国家叙事。引用多部书不自动制造全知视角；不同材料的沉默、重叠与矛盾，正是解释的一部分。

争议尤其集中在动机、数字与细节。地望、兵额、海上航程、宫廷对话和江东全境的户口都不能由一项材料裁定。较稳妥的写法是把可确认的结果与解释性的推断分开：可以说行动者在既有资源和压力下作出某种选择，也必须承认我们通常无法知道每个人在当时究竟相信什么、害怕什么。这样的保留并不会削弱历史，反而使读者看见选择为何有限、失败为何未必等于无能、胜利为何不必然证明先见。

如果把这一节点接回全章，它所显示的是一条反馈回路：舟师、港口、沿江城戍、州郡簿籍和能够随水路移动的粮食使政权能够作出行动；行动又把压力送回禁军、辅政、太子网络、州郡与水陆指挥的连接和船只、粮米、劳役和地方登记使江面上的战略最终落到家庭与县廷；新形成的权力关系反过来决定下一次资源由谁调用、沿哪条路线流动。三国史最值得阅读之处，正在于这种循环不断把宏大的名号压回具体的河流、山口、港口、文书与人事之中。诸葛恪围攻合肥新城久攻不下，吴军因疫病、补给与魏守备而撤退不是孤立条目，而是这条回路在建兴二年（253年）留下的一处可见折痕。


\subsection{从“孙峻在建业杀诸葛恪并接掌禁军与朝政”看见行动的边界}
\label{reading:tk-0253-sun-jun-coup}

建兴二年（253年），账本所记的\eventanchor{TK-0253-SUN-JUN-COUP}{孙峻在建业杀诸葛恪并接掌禁军与朝政}，首先不是抽象地图上移动的一枚棋子。读者应先把目光放在建业、武昌、濡须与长江港口：在这里，一道命令若要抵达，就必须穿过长江上下游、淮南接口、荆州峡口和长沙文书所见的地方。孙峻在建业杀诸葛恪并接掌禁军与朝政。这一结果给出可核验的时间骨架；它同时提出更难的问题：谁能调动物资与人手，谁只能等待别人的决定，谁又在道路、水面或官署的缝隙里重新计算自己的位置。

从结构上说，舟师、港口、沿江城戍、州郡簿籍和能够随水路移动的粮食并不是储存在国库里便可自动使用的力量。它们需要在不同地点被登记、运出、交接、守护和补充；其中任何一环延误，远方的军队或使节就会面对与中枢不同的现实。诸葛恪的合肥失败激化朝臣与宗室对其专权的恐惧，孙峻控制宫廷武力并利用这些不满。这段前因不应被理解为单一人物预先设计的剧本，而是一组已经存在的条件：空间给出距离，制度分配发令权，地方社会承担把命令做成行动的成本。

行动者面对的选择因而很具体。对于这一政权或集团而言，关键不是笼统地“进”或“退”，而是把下游中枢、上游防线与继承安排维持在同一条可运转的江面上。若冒险推进，可能赢得一处节点、一个名义或一次短期主动；若保留力量，又可能把接口让给竞争者，或让既有同盟与部属怀疑中枢的能力。水军维持、岸防、疾病、远征补给和宫廷清洗互相放大不是事后的修辞，它在当时已构成选择的价格，只是材料很少留下所有人的意见与犹疑。

在这一限制下，孙峻在建业杀诸葛恪并接掌禁军与朝政。把事实写到这里仍须克制：正史可以确认的通常是任命、出军、围守、退还、死亡、册封或政权转换等骨架；营中每一段传话、每一份账目和每一处战阵，往往没有同样坚实的记录。真正有解释力的，并非补足一个好看的传奇，而是观察这一骨架怎样被长江上下游、淮南接口、荆州峡口和长沙文书所见的地方的摩擦推向已经发生的结果。

直接效果首先落在可见的节点上。孙吴辅政由外戚重臣转到宗室将领孙峻手中，孙亮的皇权更加依赖禁军统帅而非公开朝议。一个城、一个渡口、一个官署、一名辅政者或一条通道的状态改变，会立刻重排相邻行动者的风险：前线可能得到喘息，近畿可能获得声望，地方官却要重新面对征发、守备或接收的命令。事件的“结束”因此只是中央记事的一条句号；在运输、任官和地方关系中，它刚刚开始产生后果。

中期影响并不等于结果的线性放大。成功会使中枢相信某条路线、某位将领或某种制度安排值得继续投资，也可能把资源更加集中到已经拥挤的方向；失利则可能迫使它缩短战线、改写同盟或重排辅政权力。就此而言，禁军、辅政、太子网络、州郡与水陆指挥的连接既是此次行动的条件，也是它留下的制度遗产。国家能力并不是“有”或“无”，而是在每一次调用之后增添某种可重复的程序，同时制造另一处更难处理的依赖。

把镜头拉近，船只、粮米、劳役和地方登记使江面上的战略最终落到家庭与县廷。史书多从朝廷、将帅和使节写起，难以连续记录沿线家庭如何应付粮食、役使、迁徙和不确定的安全；这并不意味着他们只是背景。恰恰是这些无法被完整统计的劳动，使一条道路、一支舟师或一座城戍能够在若干季节内维持。读者不宜凭缺失的户口表虚构精确的数字，却应把这种材料缺口本身看作国家抵达社会并不均匀的证据。

这一事件也改变了与他者相处的方式。把下游中枢、上游防线与继承安排维持在同一条可运转的江面上在边境上经常意味着对方也必须改算自己的粮道、使节、驻军与继承风险。即使双方暂时不战，册命、招纳、投降或撤军也会传递有关实力与脆弱性的信息。外交在这里不是战争之外的礼节，而是以较低成本测试关系的一种行动；战争则是在谈判、承诺与防守网被证明不足之后，迫使对手承担更高成本的手段。

制度所得必须和制度损失一起衡量。此次行动若使命令链更清楚、边防节点更稳固、任官更能配合资源，它就扩大了中枢的可用能力；但若它迫使权力过度依赖少数近侍、边将、转输官或地方合作者，也会让下一次继承、叛乱或远征更难收束。水军维持、岸防、疾病、远征补给和宫廷清洗互相放大由此不只是军费或伤亡的代名词，更是信息集中、地方弹性与政治信任被逐步消耗的过程。

材料在这里决定叙述能走多远。本段的基本年序和行动依据是《三国志·诸葛恪传》；《三国志·孙峻传》；《三国志·孙亮传》；《资治通鉴》魏纪。它们能够相互校验一些结果，却不是互不相干的现场档案：纪传有编纂目的，后出的通鉴重组了时间，地方文献与考古又只照亮有限地点。因此，《三国志》吴书保存宫廷主线，走马楼吴简和区域考古仅能从局部校正国家叙事。引用多部书不自动制造全知视角；不同材料的沉默、重叠与矛盾，正是解释的一部分。

争议尤其集中在动机、数字与细节。地望、兵额、海上航程、宫廷对话和江东全境的户口都不能由一项材料裁定。较稳妥的写法是把可确认的结果与解释性的推断分开：可以说行动者在既有资源和压力下作出某种选择，也必须承认我们通常无法知道每个人在当时究竟相信什么、害怕什么。这样的保留并不会削弱历史，反而使读者看见选择为何有限、失败为何未必等于无能、胜利为何不必然证明先见。

如果把这一节点接回全章，它所显示的是一条反馈回路：舟师、港口、沿江城戍、州郡簿籍和能够随水路移动的粮食使政权能够作出行动；行动又把压力送回禁军、辅政、太子网络、州郡与水陆指挥的连接和船只、粮米、劳役和地方登记使江面上的战略最终落到家庭与县廷；新形成的权力关系反过来决定下一次资源由谁调用、沿哪条路线流动。三国史最值得阅读之处，正在于这种循环不断把宏大的名号压回具体的河流、山口、港口、文书与人事之中。孙峻在建业杀诸葛恪并接掌禁军与朝政不是孤立条目，而是这条回路在建兴二年（253年）留下的一处可见折痕。


\subsection{建业、武昌、濡须与长江港口：一次转折的可达与不可达}
\label{reading:tk-0256-sun-jun-death}

太平元年（256年），账本所记的\eventanchor{TK-0256-SUN-JUN-DEATH}{孙峻去世，孙綝继掌禁军和朝政，孙亮试图摆脱权臣的空间更小}，首先不是抽象地图上移动的一枚棋子。读者应先把目光放在建业、武昌、濡须与长江港口：在这里，一道命令若要抵达，就必须穿过长江上下游、淮南接口、荆州峡口和长沙文书所见的地方。孙峻去世，孙綝继掌禁军和朝政，孙亮试图摆脱权臣的空间更小。这一结果给出可核验的时间骨架；它同时提出更难的问题：谁能调动物资与人手，谁只能等待别人的决定，谁又在道路、水面或官署的缝隙里重新计算自己的位置。

从结构上说，舟师、港口、沿江城戍、州郡簿籍和能够随水路移动的粮食并不是储存在国库里便可自动使用的力量。它们需要在不同地点被登记、运出、交接、守护和补充；其中任何一环延误，远方的军队或使节就会面对与中枢不同的现实。孙峻以宗室身份控制幼主政权后未建立可受朝廷监督的辅政机制，其死亡使军权向同族孙綝转移。这段前因不应被理解为单一人物预先设计的剧本，而是一组已经存在的条件：空间给出距离，制度分配发令权，地方社会承担把命令做成行动的成本。

行动者面对的选择因而很具体。对于这一政权或集团而言，关键不是笼统地“进”或“退”，而是把下游中枢、上游防线与继承安排维持在同一条可运转的江面上。若冒险推进，可能赢得一处节点、一个名义或一次短期主动；若保留力量，又可能把接口让给竞争者，或让既有同盟与部属怀疑中枢的能力。水军维持、岸防、疾病、远征补给和宫廷清洗互相放大不是事后的修辞，它在当时已构成选择的价格，只是材料很少留下所有人的意见与犹疑。

在这一限制下，孙峻去世，孙綝继掌禁军和朝政，孙亮试图摆脱权臣的空间更小。把事实写到这里仍须克制：正史可以确认的通常是任命、出军、围守、退还、死亡、册封或政权转换等骨架；营中每一段传话、每一份账目和每一处战阵，往往没有同样坚实的记录。真正有解释力的，并非补足一个好看的传奇，而是观察这一骨架怎样被长江上下游、淮南接口、荆州峡口和长沙文书所见的地方的摩擦推向已经发生的结果。

直接效果首先落在可见的节点上。孙綝继续支配诏令和宫禁，孙吴的继承问题与权臣控制更深地交织；宗室称谓不应掩盖具体军政网络。一个城、一个渡口、一个官署、一名辅政者或一条通道的状态改变，会立刻重排相邻行动者的风险：前线可能得到喘息，近畿可能获得声望，地方官却要重新面对征发、守备或接收的命令。事件的“结束”因此只是中央记事的一条句号；在运输、任官和地方关系中，它刚刚开始产生后果。

中期影响并不等于结果的线性放大。成功会使中枢相信某条路线、某位将领或某种制度安排值得继续投资，也可能把资源更加集中到已经拥挤的方向；失利则可能迫使它缩短战线、改写同盟或重排辅政权力。就此而言，禁军、辅政、太子网络、州郡与水陆指挥的连接既是此次行动的条件，也是它留下的制度遗产。国家能力并不是“有”或“无”，而是在每一次调用之后增添某种可重复的程序，同时制造另一处更难处理的依赖。

把镜头拉近，船只、粮米、劳役和地方登记使江面上的战略最终落到家庭与县廷。史书多从朝廷、将帅和使节写起，难以连续记录沿线家庭如何应付粮食、役使、迁徙和不确定的安全；这并不意味着他们只是背景。恰恰是这些无法被完整统计的劳动，使一条道路、一支舟师或一座城戍能够在若干季节内维持。读者不宜凭缺失的户口表虚构精确的数字，却应把这种材料缺口本身看作国家抵达社会并不均匀的证据。

这一事件也改变了与他者相处的方式。把下游中枢、上游防线与继承安排维持在同一条可运转的江面上在边境上经常意味着对方也必须改算自己的粮道、使节、驻军与继承风险。即使双方暂时不战，册命、招纳、投降或撤军也会传递有关实力与脆弱性的信息。外交在这里不是战争之外的礼节，而是以较低成本测试关系的一种行动；战争则是在谈判、承诺与防守网被证明不足之后，迫使对手承担更高成本的手段。

制度所得必须和制度损失一起衡量。此次行动若使命令链更清楚、边防节点更稳固、任官更能配合资源，它就扩大了中枢的可用能力；但若它迫使权力过度依赖少数近侍、边将、转输官或地方合作者，也会让下一次继承、叛乱或远征更难收束。水军维持、岸防、疾病、远征补给和宫廷清洗互相放大由此不只是军费或伤亡的代名词，更是信息集中、地方弹性与政治信任被逐步消耗的过程。

材料在这里决定叙述能走多远。本段的基本年序和行动依据是《三国志·孙峻传》；《三国志·孙綝传》；《三国志·孙亮传》。它们能够相互校验一些结果，却不是互不相干的现场档案：纪传有编纂目的，后出的通鉴重组了时间，地方文献与考古又只照亮有限地点。因此，《三国志》吴书保存宫廷主线，走马楼吴简和区域考古仅能从局部校正国家叙事。引用多部书不自动制造全知视角；不同材料的沉默、重叠与矛盾，正是解释的一部分。

争议尤其集中在动机、数字与细节。地望、兵额、海上航程、宫廷对话和江东全境的户口都不能由一项材料裁定。较稳妥的写法是把可确认的结果与解释性的推断分开：可以说行动者在既有资源和压力下作出某种选择，也必须承认我们通常无法知道每个人在当时究竟相信什么、害怕什么。这样的保留并不会削弱历史，反而使读者看见选择为何有限、失败为何未必等于无能、胜利为何不必然证明先见。

如果把这一节点接回全章，它所显示的是一条反馈回路：舟师、港口、沿江城戍、州郡簿籍和能够随水路移动的粮食使政权能够作出行动；行动又把压力送回禁军、辅政、太子网络、州郡与水陆指挥的连接和船只、粮米、劳役和地方登记使江面上的战略最终落到家庭与县廷；新形成的权力关系反过来决定下一次资源由谁调用、沿哪条路线流动。三国史最值得阅读之处，正在于这种循环不断把宏大的名号压回具体的河流、山口、港口、文书与人事之中。孙峻去世，孙綝继掌禁军和朝政，孙亮试图摆脱权臣的空间更小不是孤立条目，而是这条回路在太平元年（256年）留下的一处可见折痕。


\subsection{从“孙綝废孙亮，迎立琅邪王孙休，孙吴第三位皇帝即”看见行动的边界}
\label{reading:tk-0258-sun-liang-deposed}

太平三年（258年），账本所记的\eventref{TK-0258-SUN-LIANG-DEPOSED}{孙綝废孙亮，迎立琅邪王孙休，孙吴第三位皇帝即位}，首先不是抽象地图上移动的一枚棋子。读者应先把目光放在建业、武昌、濡须与长江港口：在这里，一道命令若要抵达，就必须穿过长江上下游、淮南接口、荆州峡口和长沙文书所见的地方。孙綝废孙亮，迎立琅邪王孙休，孙吴第三位皇帝即位。这一结果给出可核验的时间骨架；它同时提出更难的问题：谁能调动物资与人手，谁只能等待别人的决定，谁又在道路、水面或官署的缝隙里重新计算自己的位置。

从结构上说，舟师、港口、沿江城戍、州郡簿籍和能够随水路移动的粮食并不是储存在国库里便可自动使用的力量。它们需要在不同地点被登记、运出、交接、守护和补充；其中任何一环延误，远方的军队或使节就会面对与中枢不同的现实。孙亮试图联合近侍和宗室制衡孙綝但失败，孙綝以禁军优势操控宫廷并寻找更易接受的成年宗室。这段前因不应被理解为单一人物预先设计的剧本，而是一组已经存在的条件：空间给出距离，制度分配发令权，地方社会承担把命令做成行动的成本。

行动者面对的选择因而很具体。对于这一政权或集团而言，关键不是笼统地“进”或“退”，而是把下游中枢、上游防线与继承安排维持在同一条可运转的江面上。若冒险推进，可能赢得一处节点、一个名义或一次短期主动；若保留力量，又可能把接口让给竞争者，或让既有同盟与部属怀疑中枢的能力。水军维持、岸防、疾病、远征补给和宫廷清洗互相放大不是事后的修辞，它在当时已构成选择的价格，只是材料很少留下所有人的意见与犹疑。

在这一限制下，孙綝废孙亮，迎立琅邪王孙休，孙吴第三位皇帝即位。把事实写到这里仍须克制：正史可以确认的通常是任命、出军、围守、退还、死亡、册封或政权转换等骨架；营中每一段传话、每一份账目和每一处战阵，往往没有同样坚实的记录。真正有解释力的，并非补足一个好看的传奇，而是观察这一骨架怎样被长江上下游、淮南接口、荆州峡口和长沙文书所见的地方的摩擦推向已经发生的结果。

直接效果首先落在可见的节点上。孙吴皇帝再次由权臣更换，孙休随后必须设法处置孙綝；频繁废立削弱了吴对外动员的可信度。一个城、一个渡口、一个官署、一名辅政者或一条通道的状态改变，会立刻重排相邻行动者的风险：前线可能得到喘息，近畿可能获得声望，地方官却要重新面对征发、守备或接收的命令。事件的“结束”因此只是中央记事的一条句号；在运输、任官和地方关系中，它刚刚开始产生后果。

中期影响并不等于结果的线性放大。成功会使中枢相信某条路线、某位将领或某种制度安排值得继续投资，也可能把资源更加集中到已经拥挤的方向；失利则可能迫使它缩短战线、改写同盟或重排辅政权力。就此而言，禁军、辅政、太子网络、州郡与水陆指挥的连接既是此次行动的条件，也是它留下的制度遗产。国家能力并不是“有”或“无”，而是在每一次调用之后增添某种可重复的程序，同时制造另一处更难处理的依赖。

把镜头拉近，船只、粮米、劳役和地方登记使江面上的战略最终落到家庭与县廷。史书多从朝廷、将帅和使节写起，难以连续记录沿线家庭如何应付粮食、役使、迁徙和不确定的安全；这并不意味着他们只是背景。恰恰是这些无法被完整统计的劳动，使一条道路、一支舟师或一座城戍能够在若干季节内维持。读者不宜凭缺失的户口表虚构精确的数字，却应把这种材料缺口本身看作国家抵达社会并不均匀的证据。

这一事件也改变了与他者相处的方式。把下游中枢、上游防线与继承安排维持在同一条可运转的江面上在边境上经常意味着对方也必须改算自己的粮道、使节、驻军与继承风险。即使双方暂时不战，册命、招纳、投降或撤军也会传递有关实力与脆弱性的信息。外交在这里不是战争之外的礼节，而是以较低成本测试关系的一种行动；战争则是在谈判、承诺与防守网被证明不足之后，迫使对手承担更高成本的手段。

制度所得必须和制度损失一起衡量。此次行动若使命令链更清楚、边防节点更稳固、任官更能配合资源，它就扩大了中枢的可用能力；但若它迫使权力过度依赖少数近侍、边将、转输官或地方合作者，也会让下一次继承、叛乱或远征更难收束。水军维持、岸防、疾病、远征补给和宫廷清洗互相放大由此不只是军费或伤亡的代名词，更是信息集中、地方弹性与政治信任被逐步消耗的过程。

材料在这里决定叙述能走多远。本段的基本年序和行动依据是《三国志·孙亮传》；《三国志·孙休传》；《三国志·孙綝传》；《资治通鉴》魏纪。它们能够相互校验一些结果，却不是互不相干的现场档案：纪传有编纂目的，后出的通鉴重组了时间，地方文献与考古又只照亮有限地点。因此，《三国志》吴书保存宫廷主线，走马楼吴简和区域考古仅能从局部校正国家叙事。引用多部书不自动制造全知视角；不同材料的沉默、重叠与矛盾，正是解释的一部分。

争议尤其集中在动机、数字与细节。地望、兵额、海上航程、宫廷对话和江东全境的户口都不能由一项材料裁定。较稳妥的写法是把可确认的结果与解释性的推断分开：可以说行动者在既有资源和压力下作出某种选择，也必须承认我们通常无法知道每个人在当时究竟相信什么、害怕什么。这样的保留并不会削弱历史，反而使读者看见选择为何有限、失败为何未必等于无能、胜利为何不必然证明先见。

如果把这一节点接回全章，它所显示的是一条反馈回路：舟师、港口、沿江城戍、州郡簿籍和能够随水路移动的粮食使政权能够作出行动；行动又把压力送回禁军、辅政、太子网络、州郡与水陆指挥的连接和船只、粮米、劳役和地方登记使江面上的战略最终落到家庭与县廷；新形成的权力关系反过来决定下一次资源由谁调用、沿哪条路线流动。三国史最值得阅读之处，正在于这种循环不断把宏大的名号压回具体的河流、山口、港口、文书与人事之中。孙綝废孙亮，迎立琅邪王孙休，孙吴第三位皇帝即位不是孤立条目，而是这条回路在太平三年（258年）留下的一处可见折痕。

\begin{sourcenote}
这一组事件的来源路径并不相同。《三国志》吴书保存宫廷主线，走马楼吴简和区域考古仅能从局部校正国家叙事。阅读时应先区分能确认的年序和结果，与只能作为解释线索的动机、数量及局部过程；这比用一个完整无缝的故事覆盖材料间隙更能说明国家能力如何实际运作。
\end{sourcenote}


\subsection{一项命令怎样在长江上下游、淮南接口、荆州峡口和长沙文书所见的地方中变成现实}
\label{reading:tk-0264-sun-hao-accession}

元兴元年（264年），账本所记的\eventanchor{TK-0264-SUN-HAO-ACCESSION}{孙休去世后，孙皓被拥立为帝，孙吴进入末代皇帝统治}，首先不是抽象地图上移动的一枚棋子。读者应先把目光放在建业、武昌、濡须与长江港口：在这里，一道命令若要抵达，就必须穿过长江上下游、淮南接口、荆州峡口和长沙文书所见的地方。孙休去世后，孙皓被拥立为帝，孙吴进入末代皇帝统治。这一结果给出可核验的时间骨架；它同时提出更难的问题：谁能调动物资与人手，谁只能等待别人的决定，谁又在道路、水面或官署的缝隙里重新计算自己的位置。

从结构上说，舟师、港口、沿江城戍、州郡簿籍和能够随水路移动的粮食并不是储存在国库里便可自动使用的力量。它们需要在不同地点被登记、运出、交接、守护和补充；其中任何一环延误，远方的军队或使节就会面对与中枢不同的现实。孙休未能留下稳固成年继承人，孙吴在孙綝被诛后仍需以宗室继承重建宫廷秩序。这段前因不应被理解为单一人物预先设计的剧本，而是一组已经存在的条件：空间给出距离，制度分配发令权，地方社会承担把命令做成行动的成本。

行动者面对的选择因而很具体。对于这一政权或集团而言，关键不是笼统地“进”或“退”，而是把下游中枢、上游防线与继承安排维持在同一条可运转的江面上。若冒险推进，可能赢得一处节点、一个名义或一次短期主动；若保留力量，又可能把接口让给竞争者，或让既有同盟与部属怀疑中枢的能力。水军维持、岸防、疾病、远征补给和宫廷清洗互相放大不是事后的修辞，它在当时已构成选择的价格，只是材料很少留下所有人的意见与犹疑。

在这一限制下，孙休去世后，孙皓被拥立为帝，孙吴进入末代皇帝统治。把事实写到这里仍须克制：正史可以确认的通常是任命、出军、围守、退还、死亡、册封或政权转换等骨架；营中每一段传话、每一份账目和每一处战阵，往往没有同样坚实的记录。真正有解释力的，并非补足一个好看的传奇，而是观察这一骨架怎样被长江上下游、淮南接口、荆州峡口和长沙文书所见的地方的摩擦推向已经发生的结果。

直接效果首先落在可见的节点上。孙皓继位时魏已灭蜀、晋王将受禅，孙吴面对更孤立的国际环境；后世暴君叙事须与具体政令和地方材料区分。一个城、一个渡口、一个官署、一名辅政者或一条通道的状态改变，会立刻重排相邻行动者的风险：前线可能得到喘息，近畿可能获得声望，地方官却要重新面对征发、守备或接收的命令。事件的“结束”因此只是中央记事的一条句号；在运输、任官和地方关系中，它刚刚开始产生后果。

中期影响并不等于结果的线性放大。成功会使中枢相信某条路线、某位将领或某种制度安排值得继续投资，也可能把资源更加集中到已经拥挤的方向；失利则可能迫使它缩短战线、改写同盟或重排辅政权力。就此而言，禁军、辅政、太子网络、州郡与水陆指挥的连接既是此次行动的条件，也是它留下的制度遗产。国家能力并不是“有”或“无”，而是在每一次调用之后增添某种可重复的程序，同时制造另一处更难处理的依赖。

把镜头拉近，船只、粮米、劳役和地方登记使江面上的战略最终落到家庭与县廷。史书多从朝廷、将帅和使节写起，难以连续记录沿线家庭如何应付粮食、役使、迁徙和不确定的安全；这并不意味着他们只是背景。恰恰是这些无法被完整统计的劳动，使一条道路、一支舟师或一座城戍能够在若干季节内维持。读者不宜凭缺失的户口表虚构精确的数字，却应把这种材料缺口本身看作国家抵达社会并不均匀的证据。

这一事件也改变了与他者相处的方式。把下游中枢、上游防线与继承安排维持在同一条可运转的江面上在边境上经常意味着对方也必须改算自己的粮道、使节、驻军与继承风险。即使双方暂时不战，册命、招纳、投降或撤军也会传递有关实力与脆弱性的信息。外交在这里不是战争之外的礼节，而是以较低成本测试关系的一种行动；战争则是在谈判、承诺与防守网被证明不足之后，迫使对手承担更高成本的手段。

制度所得必须和制度损失一起衡量。此次行动若使命令链更清楚、边防节点更稳固、任官更能配合资源，它就扩大了中枢的可用能力；但若它迫使权力过度依赖少数近侍、边将、转输官或地方合作者，也会让下一次继承、叛乱或远征更难收束。水军维持、岸防、疾病、远征补给和宫廷清洗互相放大由此不只是军费或伤亡的代名词，更是信息集中、地方弹性与政治信任被逐步消耗的过程。

材料在这里决定叙述能走多远。本段的基本年序和行动依据是《三国志·孙皓传》；《三国志·孙休传》；《建康实录》；《资治通鉴》魏纪。它们能够相互校验一些结果，却不是互不相干的现场档案：纪传有编纂目的，后出的通鉴重组了时间，地方文献与考古又只照亮有限地点。因此，《三国志》吴书保存宫廷主线，走马楼吴简和区域考古仅能从局部校正国家叙事。引用多部书不自动制造全知视角；不同材料的沉默、重叠与矛盾，正是解释的一部分。

争议尤其集中在动机、数字与细节。地望、兵额、海上航程、宫廷对话和江东全境的户口都不能由一项材料裁定。较稳妥的写法是把可确认的结果与解释性的推断分开：可以说行动者在既有资源和压力下作出某种选择，也必须承认我们通常无法知道每个人在当时究竟相信什么、害怕什么。这样的保留并不会削弱历史，反而使读者看见选择为何有限、失败为何未必等于无能、胜利为何不必然证明先见。

如果把这一节点接回全章，它所显示的是一条反馈回路：舟师、港口、沿江城戍、州郡簿籍和能够随水路移动的粮食使政权能够作出行动；行动又把压力送回禁军、辅政、太子网络、州郡与水陆指挥的连接和船只、粮米、劳役和地方登记使江面上的战略最终落到家庭与县廷；新形成的权力关系反过来决定下一次资源由谁调用、沿哪条路线流动。三国史最值得阅读之处，正在于这种循环不断把宏大的名号压回具体的河流、山口、港口、文书与人事之中。孙休去世后，孙皓被拥立为帝，孙吴进入末代皇帝统治不是孤立条目，而是这条回路在元兴元年（264年）留下的一处可见折痕。


\subsection{从“西陵督步阐降晋，陆抗围西陵并击退晋援，孙吴收”看见行动的边界}
\label{reading:tk-0272-xiling}

凤凰元年（272年），账本所记的\eventanchor{TK-0272-XILING}{西陵督步阐降晋，陆抗围西陵并击退晋援，孙吴收复峡口要地}，首先不是抽象地图上移动的一枚棋子。读者应先把目光放在建业、武昌、濡须与长江港口：在这里，一道命令若要抵达，就必须穿过长江上下游、淮南接口、荆州峡口和长沙文书所见的地方。西陵督步阐降晋，陆抗围西陵并击退晋援，孙吴收复峡口要地。这一结果给出可核验的时间骨架；它同时提出更难的问题：谁能调动物资与人手，谁只能等待别人的决定，谁又在道路、水面或官署的缝隙里重新计算自己的位置。

从结构上说，舟师、港口、沿江城戍、州郡簿籍和能够随水路移动的粮食并不是储存在国库里便可自动使用的力量。它们需要在不同地点被登记、运出、交接、守护和补充；其中任何一环延误，远方的军队或使节就会面对与中枢不同的现实。步阐与孙吴朝廷关系恶化，西晋希望利用内应取得长江上游突破口，陆抗则须守住武昌以西的水陆枢纽。这段前因不应被理解为单一人物预先设计的剧本，而是一组已经存在的条件：空间给出距离，制度分配发令权，地方社会承担把命令做成行动的成本。

行动者面对的选择因而很具体。对于这一政权或集团而言，关键不是笼统地“进”或“退”，而是把下游中枢、上游防线与继承安排维持在同一条可运转的江面上。若冒险推进，可能赢得一处节点、一个名义或一次短期主动；若保留力量，又可能把接口让给竞争者，或让既有同盟与部属怀疑中枢的能力。水军维持、岸防、疾病、远征补给和宫廷清洗互相放大不是事后的修辞，它在当时已构成选择的价格，只是材料很少留下所有人的意见与犹疑。

在这一限制下，西陵督步阐降晋，陆抗围西陵并击退晋援，孙吴收复峡口要地。把事实写到这里仍须克制：正史可以确认的通常是任命、出军、围守、退还、死亡、册封或政权转换等骨架；营中每一段传话、每一份账目和每一处战阵，往往没有同样坚实的记录。真正有解释力的，并非补足一个好看的传奇，而是观察这一骨架怎样被长江上下游、淮南接口、荆州峡口和长沙文书所见的地方的摩擦推向已经发生的结果。

直接效果首先落在可见的节点上。孙吴暂时阻止晋军楔入三峡—江汉通道，但西陵危机暴露吴内部将领与中央的脆弱联系，晋继续修订伐吴方案。一个城、一个渡口、一个官署、一名辅政者或一条通道的状态改变，会立刻重排相邻行动者的风险：前线可能得到喘息，近畿可能获得声望，地方官却要重新面对征发、守备或接收的命令。事件的“结束”因此只是中央记事的一条句号；在运输、任官和地方关系中，它刚刚开始产生后果。

中期影响并不等于结果的线性放大。成功会使中枢相信某条路线、某位将领或某种制度安排值得继续投资，也可能把资源更加集中到已经拥挤的方向；失利则可能迫使它缩短战线、改写同盟或重排辅政权力。就此而言，禁军、辅政、太子网络、州郡与水陆指挥的连接既是此次行动的条件，也是它留下的制度遗产。国家能力并不是“有”或“无”，而是在每一次调用之后增添某种可重复的程序，同时制造另一处更难处理的依赖。

把镜头拉近，船只、粮米、劳役和地方登记使江面上的战略最终落到家庭与县廷。史书多从朝廷、将帅和使节写起，难以连续记录沿线家庭如何应付粮食、役使、迁徙和不确定的安全；这并不意味着他们只是背景。恰恰是这些无法被完整统计的劳动，使一条道路、一支舟师或一座城戍能够在若干季节内维持。读者不宜凭缺失的户口表虚构精确的数字，却应把这种材料缺口本身看作国家抵达社会并不均匀的证据。

这一事件也改变了与他者相处的方式。把下游中枢、上游防线与继承安排维持在同一条可运转的江面上在边境上经常意味着对方也必须改算自己的粮道、使节、驻军与继承风险。即使双方暂时不战，册命、招纳、投降或撤军也会传递有关实力与脆弱性的信息。外交在这里不是战争之外的礼节，而是以较低成本测试关系的一种行动；战争则是在谈判、承诺与防守网被证明不足之后，迫使对手承担更高成本的手段。

制度所得必须和制度损失一起衡量。此次行动若使命令链更清楚、边防节点更稳固、任官更能配合资源，它就扩大了中枢的可用能力；但若它迫使权力过度依赖少数近侍、边将、转输官或地方合作者，也会让下一次继承、叛乱或远征更难收束。水军维持、岸防、疾病、远征补给和宫廷清洗互相放大由此不只是军费或伤亡的代名词，更是信息集中、地方弹性与政治信任被逐步消耗的过程。

材料在这里决定叙述能走多远。本段的基本年序和行动依据是《三国志·陆抗传》；《晋书·羊祜传》；《晋书·杜预传》；《资治通鉴》晋纪。它们能够相互校验一些结果，却不是互不相干的现场档案：纪传有编纂目的，后出的通鉴重组了时间，地方文献与考古又只照亮有限地点。因此，《三国志》吴书保存宫廷主线，走马楼吴简和区域考古仅能从局部校正国家叙事。引用多部书不自动制造全知视角；不同材料的沉默、重叠与矛盾，正是解释的一部分。

争议尤其集中在动机、数字与细节。地望、兵额、海上航程、宫廷对话和江东全境的户口都不能由一项材料裁定。较稳妥的写法是把可确认的结果与解释性的推断分开：可以说行动者在既有资源和压力下作出某种选择，也必须承认我们通常无法知道每个人在当时究竟相信什么、害怕什么。这样的保留并不会削弱历史，反而使读者看见选择为何有限、失败为何未必等于无能、胜利为何不必然证明先见。

如果把这一节点接回全章，它所显示的是一条反馈回路：舟师、港口、沿江城戍、州郡簿籍和能够随水路移动的粮食使政权能够作出行动；行动又把压力送回禁军、辅政、太子网络、州郡与水陆指挥的连接和船只、粮米、劳役和地方登记使江面上的战略最终落到家庭与县廷；新形成的权力关系反过来决定下一次资源由谁调用、沿哪条路线流动。三国史最值得阅读之处，正在于这种循环不断把宏大的名号压回具体的河流、山口、港口、文书与人事之中。西陵督步阐降晋，陆抗围西陵并击退晋援，孙吴收复峡口要地不是孤立条目，而是这条回路在凤凰元年（272年）留下的一处可见折痕。


\subsection{从“陆抗去世，孙吴失去镇守荆州西部、熟悉晋吴边境”看见行动的边界}
\label{reading:tk-0274-lu-kang-death}

凤凰三年（274年），账本所记的\eventanchor{TK-0274-LU-KANG-DEATH}{陆抗去世，孙吴失去镇守荆州西部、熟悉晋吴边境的主要统帅}，首先不是抽象地图上移动的一枚棋子。读者应先把目光放在建业、武昌、濡须与长江港口：在这里，一道命令若要抵达，就必须穿过长江上下游、淮南接口、荆州峡口和长沙文书所见的地方。陆抗去世，孙吴失去镇守荆州西部、熟悉晋吴边境的主要统帅。这一结果给出可核验的时间骨架；它同时提出更难的问题：谁能调动物资与人手，谁只能等待别人的决定，谁又在道路、水面或官署的缝隙里重新计算自己的位置。

从结构上说，舟师、港口、沿江城戍、州郡簿籍和能够随水路移动的粮食并不是储存在国库里便可自动使用的力量。它们需要在不同地点被登记、运出、交接、守护和补充；其中任何一环延误，远方的军队或使节就会面对与中枢不同的现实。陆抗在西陵之役后维系荆州防线，其军政经验对孙吴抵御晋的上游压力至关重要。这段前因不应被理解为单一人物预先设计的剧本，而是一组已经存在的条件：空间给出距离，制度分配发令权，地方社会承担把命令做成行动的成本。

行动者面对的选择因而很具体。对于这一政权或集团而言，关键不是笼统地“进”或“退”，而是把下游中枢、上游防线与继承安排维持在同一条可运转的江面上。若冒险推进，可能赢得一处节点、一个名义或一次短期主动；若保留力量，又可能把接口让给竞争者，或让既有同盟与部属怀疑中枢的能力。水军维持、岸防、疾病、远征补给和宫廷清洗互相放大不是事后的修辞，它在当时已构成选择的价格，只是材料很少留下所有人的意见与犹疑。

在这一限制下，陆抗去世，孙吴失去镇守荆州西部、熟悉晋吴边境的主要统帅。把事实写到这里仍须克制：正史可以确认的通常是任命、出军、围守、退还、死亡、册封或政权转换等骨架；营中每一段传话、每一份账目和每一处战阵，往往没有同样坚实的记录。真正有解释力的，并非补足一个好看的传奇，而是观察这一骨架怎样被长江上下游、淮南接口、荆州峡口和长沙文书所见的地方的摩擦推向已经发生的结果。

直接效果首先落在可见的节点上。孙吴难以用一名统帅完整替代其边境网络，西部防线的协调能力下降；遗表中的战略主张须注意文本传承和追述成分。一个城、一个渡口、一个官署、一名辅政者或一条通道的状态改变，会立刻重排相邻行动者的风险：前线可能得到喘息，近畿可能获得声望，地方官却要重新面对征发、守备或接收的命令。事件的“结束”因此只是中央记事的一条句号；在运输、任官和地方关系中，它刚刚开始产生后果。

中期影响并不等于结果的线性放大。成功会使中枢相信某条路线、某位将领或某种制度安排值得继续投资，也可能把资源更加集中到已经拥挤的方向；失利则可能迫使它缩短战线、改写同盟或重排辅政权力。就此而言，禁军、辅政、太子网络、州郡与水陆指挥的连接既是此次行动的条件，也是它留下的制度遗产。国家能力并不是“有”或“无”，而是在每一次调用之后增添某种可重复的程序，同时制造另一处更难处理的依赖。

把镜头拉近，船只、粮米、劳役和地方登记使江面上的战略最终落到家庭与县廷。史书多从朝廷、将帅和使节写起，难以连续记录沿线家庭如何应付粮食、役使、迁徙和不确定的安全；这并不意味着他们只是背景。恰恰是这些无法被完整统计的劳动，使一条道路、一支舟师或一座城戍能够在若干季节内维持。读者不宜凭缺失的户口表虚构精确的数字，却应把这种材料缺口本身看作国家抵达社会并不均匀的证据。

这一事件也改变了与他者相处的方式。把下游中枢、上游防线与继承安排维持在同一条可运转的江面上在边境上经常意味着对方也必须改算自己的粮道、使节、驻军与继承风险。即使双方暂时不战，册命、招纳、投降或撤军也会传递有关实力与脆弱性的信息。外交在这里不是战争之外的礼节，而是以较低成本测试关系的一种行动；战争则是在谈判、承诺与防守网被证明不足之后，迫使对手承担更高成本的手段。

制度所得必须和制度损失一起衡量。此次行动若使命令链更清楚、边防节点更稳固、任官更能配合资源，它就扩大了中枢的可用能力；但若它迫使权力过度依赖少数近侍、边将、转输官或地方合作者，也会让下一次继承、叛乱或远征更难收束。水军维持、岸防、疾病、远征补给和宫廷清洗互相放大由此不只是军费或伤亡的代名词，更是信息集中、地方弹性与政治信任被逐步消耗的过程。

材料在这里决定叙述能走多远。本段的基本年序和行动依据是《三国志·陆抗传》；《建康实录》；《晋书·羊祜传》。它们能够相互校验一些结果，却不是互不相干的现场档案：纪传有编纂目的，后出的通鉴重组了时间，地方文献与考古又只照亮有限地点。因此，《三国志》吴书保存宫廷主线，走马楼吴简和区域考古仅能从局部校正国家叙事。引用多部书不自动制造全知视角；不同材料的沉默、重叠与矛盾，正是解释的一部分。

争议尤其集中在动机、数字与细节。地望、兵额、海上航程、宫廷对话和江东全境的户口都不能由一项材料裁定。较稳妥的写法是把可确认的结果与解释性的推断分开：可以说行动者在既有资源和压力下作出某种选择，也必须承认我们通常无法知道每个人在当时究竟相信什么、害怕什么。这样的保留并不会削弱历史，反而使读者看见选择为何有限、失败为何未必等于无能、胜利为何不必然证明先见。

如果把这一节点接回全章，它所显示的是一条反馈回路：舟师、港口、沿江城戍、州郡簿籍和能够随水路移动的粮食使政权能够作出行动；行动又把压力送回禁军、辅政、太子网络、州郡与水陆指挥的连接和船只、粮米、劳役和地方登记使江面上的战略最终落到家庭与县廷；新形成的权力关系反过来决定下一次资源由谁调用、沿哪条路线流动。三国史最值得阅读之处，正在于这种循环不断把宏大的名号压回具体的河流、山口、港口、文书与人事之中。陆抗去世，孙吴失去镇守荆州西部、熟悉晋吴边境的主要统帅不是孤立条目，而是这条回路在凤凰三年（274年）留下的一处可见折痕。

\begin{sourcenote}
这一组事件的来源路径并不相同。《三国志》吴书保存宫廷主线，走马楼吴简和区域考古仅能从局部校正国家叙事。阅读时应先区分能确认的年序和结果，与只能作为解释线索的动机、数量及局部过程；这比用一个完整无缝的故事覆盖材料间隙更能说明国家能力如何实际运作。
\end{sourcenote}

\subsection{本组事件的制度得失}

把这些节点连读，可以看出舟师、港口、沿江城戍、州郡簿籍和能够随水路移动的粮食所提供的不是无条件的优势，而是一种必须反复维护的可能。禁军、辅政、太子网络、州郡与水陆指挥的连接能够降低某些协调成本，却会让新的依赖和损失集中出现。对于读者而言，最重要的不是为每一次进退裁定单一责任，而是辨认哪些选择在当时仍然开放，哪些已被长江上下游、淮南接口、荆州峡口和长沙文书所见的地方、水军维持、岸防、疾病、远征补给和宫廷清洗互相放大和既有的人事关系缩窄。《三国志》吴书保存宫廷主线，走马楼吴简和区域考古仅能从局部校正国家叙事所保存的，正是这条能力与约束相互缠绕的历史脊柱；它没有保存的部分，则提醒我们不要把中枢语言误作全体社会的经验。

% MAINLINE-DEPTH-2026
\section{把连续节点读成一段国家能力}

以下扩展只处理已经进入账本的连续事件簇。它不为每条账本记录虚构一张独立场景卡，而是让读者在道路、城门、江面、官署与家庭负担之间追踪：资源怎样形成选择，选择怎样留下直接影响、制度所得和新的损失，材料又允许我们把话说到哪里。

\subsection{武昌与建业之间：孙吴怎样把江面做成中枢}

孙吴的中枢不是一座一次迁定的城市，而是一段必须双向调度的江面。建业面向下游交通与近畿行政，武昌保有上游与荆州方向的军事意义；两处的关联说明水路既能连接资源，也会拉长命令与救援的时间。

eventref{TK-0222-WEI-INVESTS-SUN-QUAN}{魏册孙权}、eventref{TK-0228-SHITING}{石亭之战}、eventref{TK-0229-SUN-QUAN-EMPEROR}{孙权称帝}、eventref{TK-0229-JIANYE-CAPITAL}{迁都建业}、eventref{TK-0230-WEI-WEN-MARITIME}{卫温诸葛直航海}和eventref{TK-0234-WU-HEFEI-ASSAULT}{合肥新城撤军}构成孙吴早期国家能力的几种试验。

魏的册封与质子要求试图以名义缩短江东的政治距离，孙权称帝则以相反方式把独立名号固定下来；二者都不能替代沿江州郡、水军与地方兵力的实际控制。

石亭的胜利表明吴军能利用江淮接口打断魏的攻势，航海求夷洲、亶洲与合肥进攻的挫折则表明水上组织能力并不等于任何方向都能持续投送。孙吴必须在防守、远航和进攻之间选择船只、季节与人力的用途。

直接效果是建业成为下游中枢，吴以帝国名号和水陆网络组织国家；中期效果是江面防守与上游、淮南方向的调度长期成为继承与辅政争夺也无法绕开的物质条件。

制度所得是港口、舟师和城戍形成可重复调用的防务网络；制度损失是网络维护昂贵且高度依赖季节、人员和节点，远征失败容易迅速转为宫廷人事的压力。

沿江家庭如何提供水手、粮食、船材与劳役，材料难以给出全景。走马楼吴简的局部文书可以提醒我们这类调用必须落到县廷，不能外推为全吴同一制度或人口统计。

《三国志·吴主传》、陆逊传、曹休传和《资治通鉴》给出主干；海岛位置、航程、军民死亡数和战场细节在裴注与后出材料间仍需严格分层。

孙吴的江防所以能长久，不是因为长江天然拒绝一切敌人，而是因为吴廷不断把水面、两座中枢、地方簿籍和沿岸城戍接成一套成本很高但能反复使用的网络。

\subsection{二宫与禁军：继承危机怎样进入江防的内部}

成年继承人孙登去世后，孙吴失去的不只是一个太子，而是一套已经围绕继承建立的预期。接下来，皇子宾客、官僚集团与近侍力量开始把原本属于家庭的继承问题接到国家军政的入口。

eventref{TK-0241-SUN-DENG-DEATH}{孙登去世}、eventref{TK-0242-SUN-HE-SUN-BA}{孙和与孙霸并立}、eventref{TK-0245-LU-XUN-DEATH}{陆逊去世}、eventref{TK-0250-SUN-WU-SUCCESSION-CRISIS}{二宫之争收束}、eventref{TK-0252-SUN-QUAN-DEATH}{孙权去世}、eventref{TK-0253-SUN-JUN-COUP}{孙峻政变}和eventref{TK-0258-SUN-LIANG-DEPOSED}{孙亮被废}，使继承与禁军逐步互相锁定。

两位成年皇子各有宾客与官僚支持，并不自动表示全国分裂；它却会让任官、谏争和边防判断带上继承含义。陆逊与孙权的紧张关系之所以重要，正在于上游统帅也被卷进中枢对继承的再定义。

废孙和、赐死孙霸而改立孙亮，是用强制清洗结束竞争的选择。它在短期内制造明确的名分，长期却使幼帝时期更依赖辅政者与掌握禁军的人；孙峻、孙綝的上升由此获得制度空间。

直接效果是皇位连续未中断，朝廷却连续经历辅政和废立；中期效果是建业的近侍与禁军能比成年宗室更直接地决定谁拥有皇帝名义。

制度所得是危机中仍能维持一套任命与军令的出口；制度损失是出口过度集中，边防与地方资源的实际需要不得不接受宫廷清洗的节奏。

传记常以人物忠奸和宫廷言辞组织这段故事，较难保存普通官吏、军士和地方社会如何理解连续废立。不能以传记的沉默制作全体江东的态度图。

《三国志》孙和传、孙霸传、陆逊传、孙亮传、孙峻传与裴注材料彼此补充，也保存不同的编纂立场；罪责与私语尤其不宜脱离这种材料环境。

继承危机并没有使孙吴的江防立即消失，但它让防务网络越来越需要通过不稳定的中枢入口调用；国家能力因而从外部走廊回到了宫门之内。

\subsection{东兴、西陵与陆抗：防守网络的韧性和继承的代价}

孙吴末期不宜被提前写成一片崩塌。东兴、合肥和西陵所显示的，是一个仍能以堤防、江面、城戍和熟悉边境的将领抵抗压力的政权；问题在于这种能力能否在继承、禁军与资源分配的震荡中持续。

eventref{TK-0252-DONGXING}{东兴守住濡须}、eventref{TK-0253-ZHUGE-KE-HEFEI}{合肥撤军}、eventref{TK-0264-SUN-HAO-ACCESSION}{孙皓即位}、eventref{TK-0272-XILING}{西陵收复}和eventref{TK-0274-LU-KANG-DEATH}{陆抗去世}连接起防守成功、远征受挫与上游经验流失的过程。

东兴能依托水陆地形降低守方成本，合肥方向的久攻却要求吴军把疾病、攻城、粮道和援军风险同时控制。两种结果并不矛盾，反而说明同一水网对防守和进攻的意义不同。

西陵危机中，步阐降晋与陆抗收复峡口迫使吴廷决定如何在上游保留防线。陆抗的知识并非抽象的个人才能，而是对城、江、对手部署与地方关系的长期积累；其去世使这种积累难以无成本移交。

直接效果是孙吴在二七二年仍能挫退晋援、守住重要峡口；中期效果是西晋更清楚必须把荆州、益州和淮南的压力合并，不能只依赖单一方向。

制度所得是上游守备能够在危机中协调军民和水陆节点；制度损失是它对少数边将和中枢支持的依赖加深，继承失序与人事更替更容易使边防知识断裂。

陆抗形象多由吴、晋双方传记保存。后出的赞誉或批评可提示边境关系，却不能替代完整的驻军账本，也不应把每位沿江居民的选择归结为同一忠诚。

《三国志·陆抗传》《晋书·羊祜传》《晋书·杜预传》和《资治通鉴》可交叉行动次序；城防细节、人数和边境日常仍须谨慎。

因此，西晋最终的胜利不证明孙吴此前没有能力，而是说明一套依赖江面节点和边将经验的防守体系，在多路协同与中枢持续失衡面前终于难以逐点承受。


% MAINLINE-SECOND-PASS-2026
\section{选择留下的中期压力}

事件簇的价值不在于把名称排得更密，而在于让读者看见一项选择结束之后，资源、制度和地方关系怎样继续作用。以下各节仍只接入账本已有节点；它们从直接结果向中期影响推进，并在材料不能覆盖的地方明确停下。

\subsection{两座中枢之间，江面怎样积累国家能力}

eventref{TK-0229-SUN-QUAN-EMPEROR}{孙权称帝}、eventref{TK-0229-JIANYE-CAPITAL}{迁都建业}与eventref{TK-0234-WU-HEFEI-ASSAULT}{合肥撤军}在编年中各有日期，在阅读中却应被看作同一段压力的不同切面。它们把行动者放到建业、武昌、港口舟师、濡须与淮南接口之间：这里的距离不只是路程，也是讯息延迟、转输折损、地方知识和政治信任的距离。

这一事件簇的结构条件是，政权必须将人、粮、文书、船只或军队从能够掌握的地点送到需要使用的地点。物资在出发时属于中枢，不等于抵达时仍能按中枢的计划发挥作用；道路、季节、守备与地方合作会不断改变可用资源的形态。正是这种摩擦，使看似同样的命令在不同走廊上产生不同速度。

孙吴要把下游行政、上游防线和对北试探放在同一套水陆调度内。这种取舍不是事后可以轻易取消的道德题。行动者往往只能在不完整的信息中判断：继续投入，可能取得一个节点、一次声望或短暂主动；保存力量，则可能把道路、盟友或解释权让给对手。历史的紧张处，正在于这些选择都带着可以预见但无法完全计算的代价。

直接结果通常最容易被正史保存：一座城被围或被收复，一位皇帝继位或被废，一支军队撤回，一份册命送出，一种制度语言被公布。这样的结果值得明确写出，却不是过程的全部。结果一旦进入官署与边地，才会要求新的任命、新的转输、新的守备与新的解释，原有问题因而往往以另一种形式回来。

中期影响尤其体现在信息的重新分配。一次成功会让某些将领、近侍、转输官或地方合作者被视为不可或缺；一次失利则可能让责任集中到少数可以被处分的人身上。这样形成的声望、猜疑与依赖，会影响下一次谁能发令、谁被信任、哪条路线被优先投资。国家能力因此不是战场之外的背景，而是在每一次调用中被加固或被磨损的关系。

制度所得需要和制度损失一起读。更清楚的命令链、更稳固的防线、更可重复的文书或更有效的同盟，能够减少某些协调成本；同一过程也可能使权力过度集中在少数入口，使地方弹性、替代人员和后方余裕逐步变薄。一个制度解决了今日的危机，未必同时解决它明日制造的依赖。

从家庭与地方的尺度看，行动不会只停在将领和诏书之间。运输、守城、服役、登记、接待使节和维持生产都要由沿线的人承担。现存材料很少允许我们为这些人计算整齐的数目或编造统一的态度；但这种缺口不应让他们退出叙事。恰恰相反，国家必须经过这些未被完整保存的劳动，才能把政治选择变成实际能力。

本段的证据路径以吴主传、陆逊传、满宠传、城址和航路研究为主。它们能帮助确定行动、顺序和若干结果，却具有不同的编纂位置：纪传倾向保存中枢与精英的说明，后出的编年重排前后关系，地方文书和考古只照亮局部空间。将它们并读的目的不是制造一份全知叙事，而是让每种材料约束其他材料能够承担的结论。

因此，动机、数字和细节必须维持不同的确定度。我们可以根据资源、地点和后果解释行动者为何面对某种选择；不能把这种解释写成对内心、密议或每一段行军的直接旁听。兵额、伤亡、航程、城内言辞和地方响应若缺少交叉材料，也应作为史书所述或研究争点，而不是用精确性替代证据。

两座中枢之间，江面怎样积累国家能力最终把读者带回一个共同问题：在这一段建业、武昌、港口舟师、濡须与淮南接口里，什么资源真的能被送达，谁拥有把它送达的入口，胜负之后又由谁承担制度继续运作的成本。能够回答这些问题，三国史就不再是一串已经知道结局的年份，而是一部不断在限制中重新分配选择的历史。

\subsection{二宫之争怎样把边防资源拉进宫门}

eventref{TK-0242-SUN-HE-SUN-BA}{孙和孙霸并立}、eventref{TK-0250-SUN-WU-SUCCESSION-CRISIS}{继承危机}与eventref{TK-0258-SUN-LIANG-DEPOSED}{孙亮被废}在编年中各有日期，在阅读中却应被看作同一段压力的不同切面。它们把行动者放到太子宾客、禁军、上游统帅与建业朝会之间：这里的距离不只是路程，也是讯息延迟、转输折损、地方知识和政治信任的距离。

这一事件簇的结构条件是，政权必须将人、粮、文书、船只或军队从能够掌握的地点送到需要使用的地点。物资在出发时属于中枢，不等于抵达时仍能按中枢的计划发挥作用；道路、季节、守备与地方合作会不断改变可用资源的形态。正是这种摩擦，使看似同样的命令在不同走廊上产生不同速度。

解决继承竞争需要明确名分，却会使幼帝和军令更加依赖掌握近侍的人。这种取舍不是事后可以轻易取消的道德题。行动者往往只能在不完整的信息中判断：继续投入，可能取得一个节点、一次声望或短暂主动；保存力量，则可能把道路、盟友或解释权让给对手。历史的紧张处，正在于这些选择都带着可以预见但无法完全计算的代价。

直接结果通常最容易被正史保存：一座城被围或被收复，一位皇帝继位或被废，一支军队撤回，一份册命送出，一种制度语言被公布。这样的结果值得明确写出，却不是过程的全部。结果一旦进入官署与边地，才会要求新的任命、新的转输、新的守备与新的解释，原有问题因而往往以另一种形式回来。

中期影响尤其体现在信息的重新分配。一次成功会让某些将领、近侍、转输官或地方合作者被视为不可或缺；一次失利则可能让责任集中到少数可以被处分的人身上。这样形成的声望、猜疑与依赖，会影响下一次谁能发令、谁被信任、哪条路线被优先投资。国家能力因此不是战场之外的背景，而是在每一次调用中被加固或被磨损的关系。

制度所得需要和制度损失一起读。更清楚的命令链、更稳固的防线、更可重复的文书或更有效的同盟，能够减少某些协调成本；同一过程也可能使权力过度集中在少数入口，使地方弹性、替代人员和后方余裕逐步变薄。一个制度解决了今日的危机，未必同时解决它明日制造的依赖。

从家庭与地方的尺度看，行动不会只停在将领和诏书之间。运输、守城、服役、登记、接待使节和维持生产都要由沿线的人承担。现存材料很少允许我们为这些人计算整齐的数目或编造统一的态度；但这种缺口不应让他们退出叙事。恰恰相反，国家必须经过这些未被完整保存的劳动，才能把政治选择变成实际能力。

本段的证据路径以孙和传、孙霸传、陆逊传、孙亮传和裴注为主。它们能帮助确定行动、顺序和若干结果，却具有不同的编纂位置：纪传倾向保存中枢与精英的说明，后出的编年重排前后关系，地方文书和考古只照亮局部空间。将它们并读的目的不是制造一份全知叙事，而是让每种材料约束其他材料能够承担的结论。

因此，动机、数字和细节必须维持不同的确定度。我们可以根据资源、地点和后果解释行动者为何面对某种选择；不能把这种解释写成对内心、密议或每一段行军的直接旁听。兵额、伤亡、航程、城内言辞和地方响应若缺少交叉材料，也应作为史书所述或研究争点，而不是用精确性替代证据。

二宫之争怎样把边防资源拉进宫门最终把读者带回一个共同问题：在这一段太子宾客、禁军、上游统帅与建业朝会里，什么资源真的能被送达，谁拥有把它送达的入口，胜负之后又由谁承担制度继续运作的成本。能够回答这些问题，三国史就不再是一串已经知道结局的年份，而是一部不断在限制中重新分配选择的历史。

\subsection{西陵说明孙吴末年仍有何种防守韧性}

eventref{TK-0252-DONGXING}{东兴之战}、eventref{TK-0272-XILING}{西陵收复}与eventref{TK-0274-LU-KANG-DEATH}{陆抗去世}在编年中各有日期，在阅读中却应被看作同一段压力的不同切面。它们把行动者放到堤防、峡口、沿江城戍、熟悉水陆的边将之间：这里的距离不只是路程，也是讯息延迟、转输折损、地方知识和政治信任的距离。

这一事件簇的结构条件是，政权必须将人、粮、文书、船只或军队从能够掌握的地点送到需要使用的地点。物资在出发时属于中枢，不等于抵达时仍能按中枢的计划发挥作用；道路、季节、守备与地方合作会不断改变可用资源的形态。正是这种摩擦，使看似同样的命令在不同走廊上产生不同速度。

吴廷可以利用既有节点守住方向，却要面对经验难以无成本交接的问题。这种取舍不是事后可以轻易取消的道德题。行动者往往只能在不完整的信息中判断：继续投入，可能取得一个节点、一次声望或短暂主动；保存力量，则可能把道路、盟友或解释权让给对手。历史的紧张处，正在于这些选择都带着可以预见但无法完全计算的代价。

直接结果通常最容易被正史保存：一座城被围或被收复，一位皇帝继位或被废，一支军队撤回，一份册命送出，一种制度语言被公布。这样的结果值得明确写出，却不是过程的全部。结果一旦进入官署与边地，才会要求新的任命、新的转输、新的守备与新的解释，原有问题因而往往以另一种形式回来。

中期影响尤其体现在信息的重新分配。一次成功会让某些将领、近侍、转输官或地方合作者被视为不可或缺；一次失利则可能让责任集中到少数可以被处分的人身上。这样形成的声望、猜疑与依赖，会影响下一次谁能发令、谁被信任、哪条路线被优先投资。国家能力因此不是战场之外的背景，而是在每一次调用中被加固或被磨损的关系。

制度所得需要和制度损失一起读。更清楚的命令链、更稳固的防线、更可重复的文书或更有效的同盟，能够减少某些协调成本；同一过程也可能使权力过度集中在少数入口，使地方弹性、替代人员和后方余裕逐步变薄。一个制度解决了今日的危机，未必同时解决它明日制造的依赖。

从家庭与地方的尺度看，行动不会只停在将领和诏书之间。运输、守城、服役、登记、接待使节和维持生产都要由沿线的人承担。现存材料很少允许我们为这些人计算整齐的数目或编造统一的态度；但这种缺口不应让他们退出叙事。恰恰相反，国家必须经过这些未被完整保存的劳动，才能把政治选择变成实际能力。

本段的证据路径以陆抗传、羊祜传、杜预传与边境考古研究为主。它们能帮助确定行动、顺序和若干结果，却具有不同的编纂位置：纪传倾向保存中枢与精英的说明，后出的编年重排前后关系，地方文书和考古只照亮局部空间。将它们并读的目的不是制造一份全知叙事，而是让每种材料约束其他材料能够承担的结论。

因此，动机、数字和细节必须维持不同的确定度。我们可以根据资源、地点和后果解释行动者为何面对某种选择；不能把这种解释写成对内心、密议或每一段行军的直接旁听。兵额、伤亡、航程、城内言辞和地方响应若缺少交叉材料，也应作为史书所述或研究争点，而不是用精确性替代证据。

西陵说明孙吴末年仍有何种防守韧性最终把读者带回一个共同问题：在这一段堤防、峡口、沿江城戍、熟悉水陆的边将里，什么资源真的能被送达，谁拥有把它送达的入口，胜负之后又由谁承担制度继续运作的成本。能够回答这些问题，三国史就不再是一串已经知道结局的年份，而是一部不断在限制中重新分配选择的历史。


% FINAL-READER-CLUSTERS-2026
\section{一段选择如何继续作用}

\subsection{建业的宫门和长江的港口：孙吴为何不能只选一个中枢}

eventref{TK-0229-JIANYE-CAPITAL}{迁都建业}、eventref{TK-0253-SUN-JUN-COUP}{孙峻政变}与eventref{TK-0272-XILING}{西陵收复}把读者带到建业、武昌、西陵、濡须与沿江港口。这里首先可见的是行动的表面：一项命令、一支军队、一封使节书、一场废立或一次降服。但事件之所以值得停留，正在于它们使原来分散的空间、人员和制度在同一时刻被迫相互回应。若只保留结果，便会把这种回应误写成结局早已注定。

水路网络要求中枢同时保持行政可达与上游军事反应；继承和禁军的震荡会使这两种需要相互挤压。这并不是人物品质的同义词。它是一组具体的条件：谁先得到消息，谁能调动运输与守备，谁有资格引用皇帝、联盟或法令的名义，谁又必须在道路与季节不允许时改写原先的方案。读者在这里看到的，是制度如何把可用资源转换为不均等的选择。

直接后果通常最清楚：某一方失去或取得节点，某一种人事关系被重排，某条防线得以保住或被迫后退。可是直接后果不能自动延长为中期优势。取得城池后要有人驻守，获得册封后要有人接受，压住政变后要有人重新使命令可信；这些后续工作往往比事件本身更少被完整保存。

从中期看，最重要的变化常是依赖关系的移动。一次成功会使某位将领、某个官署、一条道路或一群地方合作者显得不可替代；一次失败会迫使中枢把责任、资源和疑虑集中到另一些入口。这样形成的集中能提高下一次动员的速度，也可能让国家失去替代方案。三国的许多转折，正发生在这种效率和脆弱性同时增加的时刻。

空间在此不是沉默背景。建业、武昌、西陵、濡须与沿江港口各自要求不同的时间表：有的地点要等待水面、船只和靠泊，有的地点要计算山路、栈道和负重，有的地点则要控制宫门、城门与能够把消息写进官文书的人。行动者不能任选一种地理；他们只能在已经拥有或尚未拥有的走廊中安排风险。

把镜头再拉近，后方并不会因为史书多写将领而消失。征粮、服役、传递、维修、抄写、核验和接待使节，都意味着家庭、吏员和地方关系被卷入国家的选择。现存材料不允许我们为每一项负担虚构总数，也不允许把没有留下名字的人写成同样的态度；但正是这些未被完整记录的事务，使上层行动有可能持续。

制度所得应当与制度损失放在同一句判断中。命令链更短、防线更稳、官名更清楚或联盟更可预期，都是实际所得；对少数中军、近侍、边将、转输官和地方中介的依赖加深，也是实际损失。两者并不抵消，而是共同构成后来行动的起点。把这两面分开，历史便不会沦为把结果倒推为必然的故事。

本段主要依据《三国志》吴主传、孙峻传、陆抗传、《晋书》羊祜传及长江流域考古。它们能够为日期、行动与若干结果提供相互参照，却不是同一时间、同一立场的现场笔录。纪传通常保留精英的名义，编年著作重排前后，文书与物质材料仅照亮局部。把它们并列，正是为了让叙述既能说明国家为何行动，也能承认我们无法听见的部分。

争议也因而不可省略。兵额、伤亡、航程、密议、具体路线和每一处地方响应，常在材料之间有不同层次的可信度。较可靠的解释是从资源、空间与已知后果推出约束；较不可靠的写法，则是把后出的戏剧性叙述改写成所有行动者当时都已知的全景。

建业的宫门和长江的港口：孙吴为何不能只选一个中枢所留下的读者问题最终是：一项行动结束后，谁获得了更容易调用资源的入口，谁承担了新的风险，哪些地方关系被保留、被重排或被沉默？沿着这个问题继续阅读，三国的政治转折才会重新显出它们所依赖的道路、文书、劳动和制度代价。
